\documentclass[12pt, a4paper]{article}

% --- Шрифты и кодировки (ОБЯЗАТЕЛЬНО В ЭТОМ ПОРЯДКЕ) ---
\usepackage[T1, T2A]{fontenc} % Поддержка кириллических шрифтов в PDF
\usepackage[utf8]{inputenc}   % Кодировка исходного файла
\usepackage[russian]{babel}   % Языковой пакет

% --- Настройка полей (Стандарт для диссертаций) ---
\usepackage[left=3cm,right=2cm,top=2cm,bottom=2cm]{geometry}

% --- Математика и графика ---
\usepackage{graphicx}
\usepackage{amssymb}
\usepackage{amsmath}
\usepackage{amsthm}
\usepackage{csquotes}
\usepackage{indentfirst} % Красная строка в первом абзаце
\usepackage{caption}
\usepackage{float}
\usepackage{enumitem}
\usepackage{array}
\usepackage{tikz}
\usetikzlibrary{positioning,arrows.meta,shapes.geometric,fit,backgrounds,calc}

% --- Теоремные окружения ---
\newtheorem{theorem}{Теорема}[subsection]
\newtheorem{proposition}[theorem]{Предложение}
\newtheorem{corollary}[theorem]{Следствие}
\theoremstyle{definition}
\newtheorem{definition}[theorem]{Определение}
\theoremstyle{remark}
\newtheorem*{remark}{Замечание}

% --- Гиперссылки и метаданные ---
\usepackage{hyperref}
\hypersetup{
    unicode=true,             % Корректное отображение кириллицы в метаданных
    colorlinks=true,
    citecolor=black,
    linkcolor=black,
    filecolor=black,
    urlcolor=black,
    pdftitle={Разработка программного комплекса по оценке и защите RAG-систем от атак извлечения данных},
    pdfauthor={Волков А.Д.}
}

\begin{document}

% --- Титульная страница ---
\thispagestyle{empty}

\begin{center}
    \vspace*{-2cm}

    % Безопасная вставка логотипа
    \IfFileExists{msu.jpg}{
        \includegraphics[width=0.5\textwidth]{msu.jpg}
    }{
        \includegraphics[width=0.5\textwidth, height=4cm, demo]{msu.jpg} % Заглушка
    }
    \\[1em]

    {\scshape Московский государственный университет \\ имени М.В.~Ломоносова}\\
    Факультет вычислительной математики и кибернетики \\
    Кафедра информационной безопасности

    \vspace{2.5cm}

    {\Large Волков Артем Денисович}

    \vspace{1cm}

    {\Large Группа 619/2}

    \vspace{1cm}

    {\huge \bfseries
    Разработка программного комплекса по оценке и защите RAG-систем от атак извлечения данных\\}

    \vspace{1cm}

    {\LARGE {Магистерская диссертация}}
\end{center}

\vspace{2.5cm}

\begin{flushright}\large
  \textbf{Научный руководитель:}\\
  ассистент\\
  Е.А. Ильюшин
\end{flushright}

\vfill

\begin{center}
Москва, 2025
\end{center}

\newpage

% --- Оглавление ---
\pagenumbering{arabic}
\tableofcontents
\newpage

% --- Введение ---
\section{Введение}
\label{sec:introduction_thesis}

\subsection{Обзор предметной области и актуальность исследования}
\label{sec:overview_relevance}

Последние годы ознаменовались бурным развитием технологий в области обработки естественного языка (NLP), в первую очередь благодаря появлению и широкому распространению больших языковых моделей (LLM) \cite{Brown2020FewShot, Lewis2020RAG}. Эти модели продемонстрировали выдающиеся способности в понимании и генерации текста, но обладают рядом ограничений, таких как зависимость от фиксированного момента времени обучения и склонность к генерации недостоверной информации, известной как «галлюцинации» \cite{Qi2024Survey}.

Для преодоления данных недостатков и повышения достоверности ответов активно применяется архитектура Retrieval-Augmented Generation (RAG) \cite{Lewis2020RAG}. RAG-системы дополняют LLM возможностью поиска (Retrieval) актуальной информации во внешних источниках и используют найденные данные как контекст для генерации (Generation) ответа. Такой подход позволяет LLM давать более точные и основанные на фактах ответы, используя свежую информацию из внешних корпусов \cite{Shao2024Survey}. Благодаря своим преимуществам, RAG-системы находят все большее применение в критически важных приложениях, работающих с большими объёмами данных, таких как информационные системы, чат-боты и системы поддержки принятия решений.

Однако интеграция LLM с внешними базами знаний в рамках RAG-архитектуры порождает новые специфические уязвимости безопасности \cite{Qi2024Survey}. В частности, возникает риск \textit{атак извлечения} (Extraction Attacks), целью которых является несанкционированное получение конфиденциальной информации, содержащейся в корпусе документов, или принуждение системы к нежелательным действиям. Примеры таких атак включают Prompt Injection и атаки типа Backdoor, реализуемые через Data Poisoning \cite{Qi2023PI, Peng2024BackdoorRAG}.

На ранних этапах исследования, проводимых в рамках курсовой работы, была разработана и протестирована упрощённая RAG-система с легковесной LLM \texttt{distilgpt2}. Первоначальные результаты показали, что данная модель демонстрировала высокую устойчивость к некоторым сложным атакам (например, ASR Backdoor = 0\%), что могло давать искажённую картину реальной безопасности RAG-систем. Этот факт подчеркнул необходимость масштабирования экспериментального стенда до более мощных и реалистичных LLM. Переход на \texttt{Llama~3.1~8B} продемонстрировал, что та же базовая RAG-система без защит становится полностью уязвимой к Prompt Injection и Backdoor атакам (ASR 100\% для обоих типов). Применение базовых фильтров, разработанных в курсовой работе, снизило ASR лишь до 40\% для PI и 50\% для Backdoor, что свидетельствует о недостаточной эффективности этих мер против более продвинутых LLM. Это подтверждает, что для обеспечения безопасности RAG-систем в реальных условиях требуются более сложные и адаптивные защитные механизмы.

Актуальность темы исследования, таким образом, заключается в необходимости глубокого изучения механизмов SOTA-атак извлечения и разработки комплексных, адаптивных и экспериментально проверенных защитных мер для RAG-систем, использующих современные мощные LLM. На фоне растущего внедрения данных систем в практические приложения, особенно в сферах, где обрабатываются чувствительные или критически важные данные, обеспечение их безопасности является приоритетной задачей.

\subsection{Обзор литературы}
\label{sec:literature_review}

Архитектура Retrieval-Augmented Generation (RAG) является одним из наиболее перспективных направлений развития больших языковых моделей, позволяя преодолевать ограничения, связанные с «замороженными» знаниями и «галлюцинациями» LLM \cite{Lewis2020RAG}. Типичная RAG-система состоит из двух основных компонентов: Retriever (модуль поиска) и Generator (модуль генерации) \cite{Shao2024Survey}.

\textbf{Retriever} отвечает за поиск наиболее релевантных фрагментов информации (чанков) из внешнего корпуса документов. Этот процесс включает индексацию корпуса (разбиение документов на чанки, их векторизация с помощью эмбеддингов, например, Sentence-BERT \cite{Reimers2019SBERT}, и хранение в векторной базе данных), обработку пользовательского запроса (также векторизацию) и поиск ближайших соседей (наиболее релевантных чанков) в векторной базе \cite{Shao2024Survey}.

\textbf{Generator} (обычно LLM) получает на вход исходный запрос пользователя и найденные Retriever-ом релевантные чанки, которые формируют «контекст». LLM затем синтезирует связный и информативный ответ, основываясь как на своих внутренних знаниях, так и, преимущественно, на предоставленном внешнем контексте. Такая архитектура значительно снижает вероятность галлюцинаций и повышает отслеживаемость источника информации \cite{Lewis2020RAG}. Жизненный цикл RAG-запроса включает оффлайн-этап подготовки корпуса и онлайн-этапы получения, векторизации запроса, поиска чанков, формирования контекста, генерации и выдачи ответа пользователю \cite{Shao2024Survey}.

Несмотря на преимущества, RAG-системы наследуют общие риски LLM и приобретают специфические уязвимости \cite{Qi2024Survey}. Общие риски LLM включают классический Prompt Injection (манипуляция поведением LLM через вредоносные инструкции в запросе), галлюцинации и предвзятость \cite{Qi2024Survey}.

Специфические риски RAG-систем возникают из-за зависимости от внешнего корпуса данных и процесса извлечения:
\begin{itemize}[label=$\bullet$]
    \item \textbf{Извлечение конфиденциальной информации}: основная цель атак извлечения — получить несанкционированный доступ к чувствительной информации из корпуса документов. Это возможно как через Prompt Injection, так и через Data Poisoning, заставляя LLM выдать конфиденциальные данные \cite{Qi2024Survey, Peng2024BackdoorRAG}.
    \item \textbf{Манипуляция релевантностью}: атакующий может попытаться повлиять на процесс ретривала, чтобы обеспечить включение вредоносных чанков или игнорирование важных \cite{Peng2024BackdoorRAG}.
\end{itemize}

Ключевыми механизмами атак извлечения являются:
\begin{itemize}[label=$\bullet$]
    \item \textbf{Prompt Injection (PI)}: атакующий конструирует запрос, содержащий легитимную часть и вредоносные инструкции. RAG-система извлекает релевантный (потенциально чувствительный) контекст. LLM получает полный промпт, и из-за конкуренции инструкций может отдать приоритет вредоносным указаниям, игнорируя системные правила, что приводит к извлечению информации \cite{Qi2023PI}.
    \item \textbf{Атаки через данные (Data Poisoning, Backdoor)}: злоумышленник создаёт «отравленный» чанк с вредоносной инструкцией (payload) и триггерной фразой, скрытно внедряет его в корпус документов. При получении запроса с триггерной фразой Retriever извлекает отравленный чанк, и LLM выполняет вредоносную инструкцию, раскрывая конфиденциальную информацию \cite{Gu2017Badnets, Peng2024BackdoorRAG}.
\end{itemize}

Современные исследования активно развивают как методы атак, так и механизмы защиты. Среди SOTA-атак выделяется \textbf{External Data Extraction Attack (EDEA)} и метод \textbf{SECRET} \cite{SECRET2025}, который использует LLM для автоматической оптимизации Jailbreak-промптов и кластерный триггеринг для эффективного поиска. Такие атаки демонстрируют значительно более высокую эффективность по сравнению с традиционными методами \cite{SECRET2025}. Важной особенностью SECRET является \textit{структурная декомпозиция} атакующего запроса на три функциональных компонента — инструкцию извлечения, jailbreak-оператор и триггер ретривала. Это свойство, как будет показано в разделе~\ref{sec:def_hard}, является не только характеристикой атаки, но и точкой опоры для её детектирования.

В ответ на эти угрозы разрабатываются новые парадигмы защиты:
\begin{itemize}[label=$\bullet$]
    \item \textbf{LeakSealer} \cite{LeakSealer2025}: предлагает поведенческий анализ, комбинируя статический анализ исторических запросов с динамической защитой для определения аномалий путём группировки запросов по темам и их классификации с помощью LLM.
    \item \textbf{ControlNET} \cite{ControlNET2025}: функционирует как «файрвол» на основе анализа внутренних активаций LLM. Он анализирует входящие и исходящие данные, используя технические сигналы, такие как \textit{activation shift}, чтобы обнаружить семантически подозрительные запросы или нежелательное «мышление» LLM.
    \item Архитектура \textbf{RAGFort} \cite{RAGFort2025}: представляет собой двухуровневую защиту, нацеленную на атаки внутри одной темы (\textit{intra-class}) и с переходом между темами (\textit{inter-class}). Она включает \textbf{Contrastive Re-indexing} для семантической изоляции разных тем в индексе и \textbf{Constrained Cascade Generation} — каскадную генерацию ответа с контролем по оценке неопределённости генератора.
\end{itemize}
Эти подходы демонстрируют, что эффективная защита должна быть комплексной, адаптивной и работать на нескольких уровнях архитектуры RAG. Однако каждая из перечисленных защит по отдельности работает лишь с одним каналом сигнала (эмбеддинг запроса, активации LLM, ответ генератора соответственно) и не использует структуру современной SOTA-атаки SECRET, что открывает возможность для построения комбинированной защиты с повышенной устойчивостью.

\subsection{Постановка задачи исследования}
\label{sec:problem_statement}

С учётом возрастающей популярности и использования RAG-систем в приложениях, обрабатывающих конфиденциальные данные, а также недостаточной эффективности базовых защитных мер против мощных современных LLM (как показал переход на \texttt{Llama~3.1}), возрастает критическая необходимость в разработке и проверке более сложных систем безопасности. Существующие SOTA-атаки (например, SECRET) представляют серьёзную угрозу для целостности, конфиденциальности и доступности информации, требуя разработки адекватных адаптивных защитных механизмов.

Целью данного исследования является \textbf{разработка программного комплекса} для оценки и защиты RAG-систем от атак извлечения данных. Программный комплекс должен поддерживать сборку RAG-систем из взаимозаменяемых компонентов (векторное хранилище, эмбеддер, LLM-провайдер), моделирование актуальных атак извлечения, подключение защитных профилей, измерение метрик качества и устойчивости и предоставление пользователю интерфейса для проведения экспериментов. Алгоритмическим ядром защитной составляющей комплекса выступает оригинальный каскадный адаптивный метод HARD, разработанный в рамках настоящей работы.

Предлагаемый метод HARD (\textit{Hybrid Adaptive RAG Defense}) строится как каскадная адаптация подхода RAGFort (\textit{Constrained Cascade Generation}, \textit{Contrastive Re-indexing}) к сценарию black-box доступа к LLM (см.~\S\ref{sec:threat_model}) и дополняется двумя оригинальными компонентами:
\begin{itemize}[label=$\bullet$]
    \item \textbf{Intent--Retrieval Dissociation (IRD)} — детектор атак, эксплуатирующий структурную декомпозируемость SOTA-атак SECRET (см. формулу \eqref{eq:secret_decomp}) через анализ расхождения ретривальных отпечатков семантических фрагментов запроса;
    \item \textbf{Risk-Budget Thresholding} — каскадное распространение риск-бюджета, делающее пороги последующих стадий защиты функцией агрегированного риска предыдущих стадий.
\end{itemize}

Для обеспечения строгости и ясности исследования основные сущности и процессы предварительно формализованы следующим образом (подробная формализация представлена в главе~\ref{sec:math_formalization_chapter}):
\begin{itemize}[label=$\square$]
    \item \textbf{RAG-система} может быть представлена как композиция функций: $f_{RAG}(q, D) = G(q, R(q, D))$, где $q$~--- пользовательский запрос, $D$~--- корпус документов, $R(\cdot)$~--- функция Retriever'а, извлекающая контекст $C \subset D$, и $G(\cdot)$~--- функция Generator'а (LLM), формирующая ответ на основе запроса и контекста.
    \item \textbf{Атака извлечения данных}~--- это модификация запроса $q \to q_{atk}$ и/или корпуса $D \to D_{atk}$ с целью заставить RAG-систему выдать конфиденциальную информацию $S$. Атака считается успешной, если $f_{extract}(f_{RAG}(q_{atk}, D_{atk})) = 1$, где $f_{extract}(\cdot)$~--- детектор утечки информации $S$.
    \item \textbf{Защитный механизм}~--- это функция $Def(\cdot)$, которая модифицирует поведение RAG-системы $f_{RAG}$ до $f'_{RAG}$ таким образом, чтобы минимизировать успешность атак: $f'_{RAG}(q, D) = Def(f_{RAG}(q, D))$, стремясь к $f_{extract}(f'_{RAG}(q_{atk}, D_{atk})) \to 0$ при сохранении функциональности для легитимных запросов.
\end{itemize}

\subsection{Задачи исследования}
\label{sec:research_objectives}

Для достижения поставленной цели были сформулированы следующие задачи:
\begin{enumerate}[label=\arabic*.]
    \item Провести всестороннее исследование архитектуры RAG-систем, специфических угроз безопасности (Prompt Injection и Backdoor/Data Poisoning), а также детально проанализировать актуальные SOTA-атаки (EDEA и метод SECRET) и SOTA-защиты (LeakSealer, ControlNET, RAGFort), включая новые исследования 2024--2025 гг.
    \item Построить единый формальный каркас каскадной защиты, сформулировать и доказать границу вероятности успешной атаки для каскадного защитного оператора (\textit{Cascade Bound}) при условии диверсификации сигналов стадий.
    \item Разработать метод \textbf{HARD} и его два оригинальных компонента --- сигнал \textbf{IRD} на основе структурной декомпозиции SECRET-атак и механизм \textbf{Risk-Budget Thresholding} адаптации порогов.
    \item Разработать программный комплекс для воспроизводимой оценки и защиты RAG-систем, поддерживающий: (а)~сборку RAG-системы из взаимозаменяемых компонентов (векторное хранилище, эмбеддер, LLM по ролям); (б)~подключение защитных профилей по имени; (в)~моделирование классов атак извлечения; (г)~автоматический сбор метрик и формирование отчётов; (д)~интерактивный пользовательский интерфейс для запуска экспериментов и выбора компонентов без правки исходного кода.
    \item Реализовать и протестировать SOTA-атаку типа SECRET на разработанном стенде с \texttt{Llama~3.1}, измерить её базовую успешность (ASR) и проверить устойчивость текущих базовых фильтров и прототипа LLM-Guard.
    \item Реализовать метод HARD в составе программного комплекса и провести всестороннюю экспериментальную оценку его эффективности против SOTA-атак (SECRET), Prompt Injection и Backdoor, а также на легитимных запросах, используя метрики ASR, FPR, BPD, Latency Overhead, а также вспомогательные метрики StagePassRate, StageCorrelation и $\mathrm{AUC}_{\text{IRD}}$ для эмпирической проверки допущений теоретической модели.
    \item Сформулировать детальные рекомендации по повышению безопасности RAG-систем и определить перспективные направления дальнейших исследований в данной области.
\end{enumerate}

\subsection{План решения}
\label{sec:solution_plan}

На текущем этапе исследования была проделана следующая работа:
\begin{itemize}[label=$\bullet$]
    \item Проведена \textbf{математическая формализация} всех методов и компонентов RAG-систем, атак извлечения и защитных механизмов, построена единая модель каскадного защитного оператора, сформулирована и доказана каскадная граница успеха атаки (теорема~\ref{thm:cascade_bound}).
    \item Предложен и формализован метод \textbf{HARD} с двумя оригинальными компонентами --- сигналом \textbf{IRD} и механизмом \textbf{Risk-Budget Thresholding}.
    \item Подтверждено, что легковесная \texttt{distilgpt2} давала искажённую картину безопасности из-за своей нечувствительности к сложным атакам (ASR Backdoor = 0\%).
    \item Произведено масштабирование экспериментального стенда до \texttt{Llama~3.1~8B}. Выявлено, что на данной LLM базовая RAG-система без защит имеет ASR 100\% для Prompt Injection и Backdoor, а базовые фильтры из курсовой работы снижают ASR лишь до 40\%/50\% соответственно.
    \item Детально изучена и \textbf{реализована атака SECRET} \cite{SECRET2025} на экспериментальном стенде с \texttt{Llama~3.1}, включая разработку Jailbreak-операторов и триггеров извлечения.
    \item Измерен базовый ASR атаки SECRET на \texttt{Llama~3.1} без защит. Протестирована устойчивость текущего \texttt{LLM-Guard} и базовых фильтров против атаки SECRET, что позволило определить текущий уровень защищённости.
    \item Разработан и протестирован прототип \texttt{LLM-Guard}, использующий LLM-«охранника» и guard-промпты для семантического анализа входных запросов, финальных промптов и выходных ответов. На \texttt{Llama~3.1} этот прототип показал ASR 80\% для PI и 0\% для Backdoor. Данный \texttt{LLM-Guard} является практическим прототипом идеи LLM-«верификатора» из архитектуры RAGFort.
    \item Проведён анализ SOTA-исследований в области атак (EDEA, SECRET \cite{SECRET2025}) и защит (LeakSealer \cite{LeakSealer2025}, ControlNET \cite{ControlNET2025}, RAGFort \cite{RAGFort2025}).
\end{itemize}

\textbf{План дальнейших работ} в рамках диссертационного исследования включает следующие этапы:

\textbf{Фаза 2: Реализация метода HARD}
\begin{itemize}[label=$\circ$]
    \item Реализовать модуль \textbf{IRD}, включая семантическую декомпозицию запроса $\phi$ на функциональные фрагменты и вычисление среднего попарного расстояния Жаккара между ретривальными отпечатками (формула \eqref{eq:ird_score}).
    \item Реализовать модуль \textbf{Topic-Consistent Re-ranking (TCR)} на базе оффлайн-кластеризации корпуса, адаптирующий идею Contrastive Re-indexing из RAGFort к онлайн-сценарию.
    \item Реализовать каскадную генерацию \textit{Draft-then-Verify} как суррогат Constrained Cascade Generation из RAGFort, пригодный для сценария black-box API доступа к LLM.
    \item Реализовать \textbf{Post-processing Guard} (лексический и $\mathrm{LCSR}$-сканер финального ответа) как четвёртую стадию каскада.
    \item Реализовать механизм \textbf{Risk-Budget Thresholding} и единый оркестратор каскада, обеспечивающий адаптивную настройку порогов.
\end{itemize}

\textbf{Фаза 3: Финальная оценка и анализ}
\begin{itemize}[label=$\circ$]
    \item Провести полный цикл комплексных тестов SOTA-атаки \textbf{SECRET}, а также Prompt Injection и Backdoor против разработанного метода HARD и его ablation-версий.
    \item Собрать финальные метрики эффективности (ASR, FPR, BPD, Latency Overhead), а также вспомогательные метрики (StagePassRate, StageCorrelation, $\mathrm{AUC}_{\text{IRD}}$, $\kappa$) для эмпирической проверки теоретической границы из теоремы~\ref{thm:cascade_bound}.
    \item Провести сравнительный анализ эффективности HARD с базовыми фильтрами (уровня курсовой работы) и с полной black-box-адаптацией RAGFort (профиль \texttt{ragfort} = Draft-then-Verify $+$ Topic-Consistent Re-ranking) как с главным baseline. Это позволяет оценить вклад именно оригинальных компонентов настоящей работы~--- IRD, Output Leak Scanner и Risk-Budget Thresholding~--- сверх RAGFort-адаптации.
    \item Подготовить окончательные выводы для магистерской диссертации, включая рекомендации по внедрению и направления дальнейших исследований.
\end{itemize}

% --- Глава 2: Математическая формализация ---
\section{Математическая формализация моделей атак и методов защиты}
\label{sec:math_formalization_chapter}

В данной главе вводится формальный математический аппарат, описывающий функционирование систем Retrieval-Augmented Generation (RAG), модели угроз, а также методы атак и защиты. Формализация необходима для строгого определения целей злоумышленника, формулировки гарантий защитных механизмов и метрик эффективности.

\subsection{Модель угроз}
\label{sec:threat_model}

Принятая в настоящей работе модель угроз~--- \textbf{строго black-box}: атакующий взаимодействует с RAG-системой только через её интерфейс ответов и не имеет доступа к внутренним состояниям LLM-генератора (активации скрытых слоёв, логиты, веса, градиенты), к файлам векторного индекса и к серверной инфраструктуре жертвы. Возможности атакующего:
\begin{itemize}[label=$\bullet$]
    \item \textbf{Запросная сторона.} Произвольно конструирует запрос $q \in \mathcal{V}^*$ и наблюдает финальный ответ $a$ системы; внутренние сигналы (вероятности, активации) недоступны.
    \item \textbf{Корпус.} Не имеет прямого доступа к индексу $\mathcal{I}$, к эмбеддеру $E$ и к файлам корпуса $\mathcal{D}$ на стороне жертвы, но способен внедрить в корпус отравленный документ $d_{pois}$ \textit{через любой штатный канал ингеста}. К таким каналам относятся, в частности: публикация контента в открытых или полуоткрытых источниках, из которых жертва автоматически собирает корпус (вики, форумы, новостные ленты, RSS); загрузка документа через пользовательский интерфейс системы, если такая возможность предоставляется (например, корпоративная база знаний, в которой произвольный сотрудник или контрагент может добавить документ); подача документа на согласование администратору, который штатно принимает поступления; внесение изменений в индексируемые репозитории документации. Конкретный канал ингеста для модели угроз и для метода защиты несущественен; существенно лишь то, что отравленный документ проходит \textit{обычный pipeline эмбеддинга} ($d_{pois} \to E(d_{pois}) \to \mathcal{I}$) и оказывается в индексе наравне с легитимными чанками. Прямая модификация уже построенного векторного индекса, эмбеддера или весов LLM выходит за пределы рассматриваемой модели угроз и относится к классу инфраструктурной компрометации, защита от которой обеспечивается стандартными средствами контроля доступа.
\end{itemize}

Защитная сторона (RAG-провайдер) также рассматривается в black-box режиме относительно LLM-генератора: предполагается, что генератор доступен только через API (типичная ситуация для коммерческих LLM-сервисов и большинства корпоративных развёртываний с внешними моделями). Это исключает из рассмотрения защиты, требующие доступа к внутренним активациям LLM, и определяет выбор практических суррогатов для тех SOTA-методов, которые в оригинале опираются на токен-уровневое стробирование или анализ скрытых состояний (см. ремарки в~\S\ref{sec:def_sota}).

\subsection{Формальная модель RAG-системы}
\label{sec:math_rag_model}

Пусть $\mathcal{V}$~--- словарь токенов. Обозначим пространство всех возможных последовательностей токенов как $\mathcal{V}^*$. Система RAG определяется кортежем $\langle \mathcal{D}, E, \mathcal{I}, R, G \rangle$, где:
\begin{itemize}
    \item $\mathcal{D} = \{d_1, d_2, \dots, d_N\}$~--- корпус документов, где каждый документ $d_i \in \mathcal{V}^*$;
    \item $E: \mathcal{V}^* \to \mathbb{R}^k$~--- функция эмбеддинга, отображающая текст в $k$-мерное векторное пространство (например, Sentence-BERT);
    \item $\mathcal{I} = \{ (\mathbf{v}_i, \text{meta}_i) \}_{i=1}^M$~--- векторный индекс, где $\mathbf{v}_i = E(c_i)$, а $c_i$~--- фрагмент (чанк) документа $d_j \in \mathcal{D}$;
    \item $R: \mathcal{V}^* \times \mathcal{I} \to \mathcal{P}(\mathcal{V}^*)$~--- функция ретривера (Retriever), возвращающая подмножество чанков $C \subset \mathcal{D}$ для запроса $q$;
    \item $G: \mathcal{V}^* \times \mathcal{V}^* \to \mathcal{V}^*$~--- генеративная модель (Generator/LLM), принимающая на вход контекст и запрос.
\end{itemize}

\subsubsection{Процесс генерации ответа}
Для пользовательского запроса $q \in \mathcal{V}^*$:

1.~\textbf{Поиск (Retrieval):} вычисляется множество релевантных чанков $C_q = \{c_1, \dots, c_k\}$, максимизирующих меру сходства $sim(\cdot, \cdot)$ (обычно косинусное сходство):
    \begin{equation}
        C_q = \underset{C \subset \mathcal{D}, |C|=k}{\arg\max} \sum_{c \in C} sim(E(q), E(c))
    \end{equation}

2.~\textbf{Формирование промпта:} контекст объединяется с запросом с использованием шаблонной функции $\mathcal{T}$:
    \begin{equation}
        p = \mathcal{T}(q, C_q) = \text{Context: } [c_1, \dots, c_k];\ \text{Query: } q
    \end{equation}

3.~\textbf{Генерация (Generation):} ответ $a$ генерируется авторегрессионно, максимизируя правдоподобие следующего токена:
\begin{equation}
    a = G(p) = \underset{y \in \mathcal{V}^*}{\arg\max} \prod_{t=1}^{|y|} P(y_t \mid y_{<t}, p; \theta_G)
\end{equation}
где
\begin{itemize}
    \item $a$~--- итоговая последовательность токенов (ответ системы);
    \item $y_t$~--- токен, генерируемый на шаге $t$;
    \item $y_{<t}$~--- последовательность токенов, сгенерированных на предыдущих шагах (контекст генерации);
    \item $P(y_t \mid y_{<t}, p; \theta_G)$~--- условная вероятность появления токена $y_t$ при условии поданного промпта $p$ и ранее сгенерированного текста, определяемая параметрами нейросети $\theta_G$.
\end{itemize}

\subsection{Математическая модель атак извлечения данных}
\label{sec:math_attacks}

Целью атаки извлечения данных является получение конфиденциальной информации $S_{priv} \subset \mathcal{D}$, которая не должна быть выдана пользователю. Обозначим функцию проверки утечки как индикаторную функцию $\mathbb{I}_{leak}(a, S_{priv})$, которая равна $1$, если ответ $a$ содержит информацию из $S_{priv}$, и $0$ иначе. Формально функцию проверки утечки можно определить как отображение $\mathbb{I}_{leak}: \mathcal{V}^* \times \mathcal{V}^* \to \{0, 1\}$:
\begin{equation}
    \mathbb{I}_{leak}(a, S_{priv}) =
    \begin{cases}
      1, & \text{если } S_{priv} \sqsubseteq a \\
      0, & \text{иначе}
    \end{cases}
\end{equation}
где $a \in \mathcal{V}^*$~--- сгенерированный ответ, $S_{priv} \in \mathcal{V}^*$~--- целевая конфиденциальная информация, а отношение $\sqsubseteq$ обозначает вхождение подстроки.

\subsubsection{Prompt Injection (PI)}
\label{sec:pi_formalization}
В атаке Prompt Injection злоумышленник не имеет доступа к корпусу $\mathcal{D}$, но контролирует запрос $q$. Запрос модифицируется в $q' = q_{benign} \oplus \delta$, где $\delta$~--- вредоносная инъекция.

Задача оптимизации для атакующего:
\begin{equation}
    \delta^* = \underset{\delta \in \mathcal{V}^*}{\arg\max} \; \mathbb{P}\Big( \mathbb{I}_{leak}\big( G(\mathcal{T}(q_{benign} \oplus \delta, R(q', \mathcal{I}))), S_{priv} \big) = 1 \Big)
\end{equation}
при ограничении на длину инъекции $|\delta| < L$ (контекстное окно + скрытность).

Где элементы формулы интерпретируются следующим образом:
\begin{itemize}
    \item $\delta^*$~--- искомая оптимальная вредоносная инъекция (строка-атака);
    \item $q_{benign}$~--- легитимная часть запроса, служащая для маскировки;
    \item $\oplus$~--- операция конкатенации строк;
    \item $R(q', \mathcal{I})$~--- результат работы ретривера: множество документов, извлечённых по модифицированному запросу $q'$;
    \item $G(\mathcal{T}(\dots))$~--- ответ, сгенерированный LLM на основе сформированного промпта;
    \item $S_{priv}$~--- конфиденциальная информация (цель атаки);
    \item $\mathbb{I}_{leak}(a, S_{priv})$~--- индикаторная функция успеха атаки;
    \item $\mathbb{P}(\dots)$~--- вероятность успешного извлечения секрета (учитывая стохастическую природу генерации LLM).
\end{itemize}

\subsubsection{Data Poisoning / Backdoor}
\label{sec:backdoor_formalization}
Злоумышленник имеет возможность внедрить в корпус $\mathcal{D}$ отравленный документ $d_{pois}$ через один из штатных каналов ингеста, перечисленных в~\S\ref{sec:threat_model}. Документ $d_{pois}$ содержит триггерную последовательность $\tau$ и вредоносную полезную нагрузку (payload) $\pi$. С точки зрения формализации защиты канал ингеста несущественен: существенно лишь то, что $d_{pois}$ проходит штатный pipeline эмбеддинга и оказывается в индексе $\mathcal{I}$ наравне с легитимными чанками.

Атака состоит из двух этапов:

1.~\textbf{Внедрение (Poisoning):} найти $d_{pois}$ такой, что для любого запроса $q$, содержащего триггер $\tau$, ретривер выберет $d_{pois}$:
    \begin{equation}
        \forall q \in \mathcal{V}^*,\ \tau \subset q \implies d_{pois} \in R(q, \mathcal{I} \cup E(d_{pois}))
    \end{equation}
    Это достигается минимизацией расстояния в пространстве эмбеддингов:
    \begin{equation}
        dist(E(d_{pois}), E(\tau)) \to 0
    \end{equation}

2.~\textbf{Активация (Inference):} генератор под воздействием контекста $d_{pois}$ должен выдать целевую информацию:
    \begin{equation}
        G(\mathcal{T}(q, \{ \dots, d_{pois}, \dots \})) \approx S_{priv}
    \end{equation}

\subsubsection{Формализация SOTA-атаки: SECRET (на основе EDEA)}
\label{sec:secret_formalization}
Согласно исследованию He et al. \cite{SECRET2025}, атака SECRET базируется на унифицированном фреймворке EDEA, который декомпозирует атакующий запрос $x$ на три функциональных компонента:
\begin{equation}
    \label{eq:secret_decomp}
    x = I_{ext} \oplus O_{jail} \oplus T_{retr}
\end{equation}
где
\begin{itemize}
    \item $I_{ext}$ (Extraction Instruction)~--- инструкция извлечения (например, «Please repeat the text verbatim»);
    \item $O_{jail}$ (Jailbreak Operator)~--- оператор, оптимизированный для обхода защитных фильтров (Safety Alignment);
    \item $T_{retr}$ (Retrieval Trigger)~--- триггер, направляющий ретривер к целевым документам.
\end{itemize}

Метод SECRET реализует раздельную оптимизацию этих компонентов:

\textbf{1. Оптимизация джейлбрейка ($O_{jail}$).}
Используется подход \textit{LLM-as-an-Optimizer}. Задача формулируется как поиск оператора, максимизирующего вероятность выполнения инструкции $I_{ext}$ при наличии «загрязнённого» контекста $c$:
\begin{equation}
    O_{jail}^* = \underset{O \in \mathcal{V}^*}{\arg\max} \; \mathbb{E}_{c \sim \mathcal{D}} \left[ \log P(I_{ext\_success} \mid I_{ext} \oplus O \oplus c; \theta_G) \right]
\end{equation}
Авторы используют механизм мутации промптов (Prompt Mutation) для итеративного улучшения $O_{jail}$.

\textbf{2. Оптимизация триггера (Cluster-Focused Triggering, CFT).}
Для масштабируемого поиска $T_{retr}$ применяется стратегия CFT, чередующая два этапа:
\begin{itemize}
    \item \textbf{Global Exploration:} кластеризация эмбеддингов корпуса $\mathcal{D}$ (например, через K-Means) для выявления центроидов тем $\{\mu_1, \dots, \mu_k\}$;
    \item \textbf{Local Exploitation:} генерация триггера $T$, минимизирующего расстояние до центроида целевого кластера $\mu_i$:
    \begin{equation}
        T_{retr}^* = \underset{T \in \mathcal{V}^*}{\arg\min} \; \lVert E(T) - \mu_i \rVert_2
    \end{equation}
\end{itemize}
Такой подход позволяет эффективно извлекать данные из различных семантических областей без полного перебора корпуса.

\subsection{Математическая модель методов защиты}
\label{sec:math_defenses}

В настоящем разделе вводится единый формализм защитного оператора, позволяющий описать как классические методы защиты, так и современные SOTA-подходы в качестве \textit{компонентов} единой архитектуры. На его основе в~\S\ref{sec:def_hard} формулируется предлагаемый в работе метод HARD.

\subsubsection{Общий формализм защитного оператора}
\label{sec:def_general}

Обозначим конфигурационное пространство RAG-системы $\mathcal{S} = \langle \mathcal{D}, E, \mathcal{I}, R, G, \mathcal{T} \rangle$, а ответ системы на запрос $q$~--- как $f_{RAG}(q; \mathcal{S}) \in \mathcal{V}^*$. Множество атакующих запросов обозначим $\mathcal{Q}_{atk}$, множество легитимных~--- $\mathcal{Q}_{legit}$; $S_{priv} \subset \mathcal{V}^*$~--- целевая конфиденциальная информация.

\begin{definition}[Защитный оператор]
\label{def:defense_operator}
\textit{Механизмом защиты} будем называть оператор $\mathrm{Def}$, отображающий пару $(q, \mathcal{S})$ в ответ системы или специальный безопасный ответ $a_\varnothing \in \mathcal{V}^*$:
\begin{equation}
    \mathrm{Def}: (q, \mathcal{S}) \mapsto a \in \mathcal{V}^* \cup \{a_\varnothing\}.
\end{equation}
\end{definition}

\begin{definition}[Задача построения эффективной защиты]
\label{def:optimal_defense}
\textit{Эффективным} будем называть механизм защиты, решающий задачу
\begin{equation}
    \label{eq:def_objective}
    \mathrm{Def}^{*} = \underset{\mathrm{Def}}{\arg\min}
    \frac{1}{|\mathcal{Q}_{atk}|}
    \sum_{q \in \mathcal{Q}_{atk}}
    \mathbb{I}_{leak}\!\bigl(f_{RAG}(q; \mathrm{Def}(\mathcal{S})),\ S_{priv}\bigr),
\end{equation}
при ограничениях $\mathrm{FPR}(\mathrm{Def}) \le \varepsilon_{fpr}$ и $\mathrm{BPD}(\mathrm{Def}) \le \varepsilon_{bpd}$, определённых на множестве $\mathcal{Q}_{legit}$ (см.~\S\ref{sec:metrics}). Ограничения исключают тривиальное решение «блокировать всё».
\end{definition}

\paragraph{Каскадная декомпозиция.}
На практике защита реализуется как композиция стадий, соответствующих этапам жизненного цикла RAG-запроса:
\begin{equation}
    \label{eq:def_cascade}
    \mathrm{Def} \;=\;
    \mathrm{Def}_{out} \circ
    \mathrm{Def}_{gen} \circ
    \mathrm{Def}_{ret} \circ
    \mathrm{Def}_{in}.
\end{equation}
Каждая стадия $\mathrm{Def}_i$~--- это тройка $\langle h_i, \theta_i, \text{act}_i \rangle$, где:
\begin{itemize}
    \item $h_i: \mathcal{X}_i \to [0, 1]$~--- \textit{сигнальная функция риска}, принимающая на вход доступные стадии наблюдения $x_i$ (запрос, контекст, черновик ответа, финальный ответ);
    \item $\theta_i \in [0, 1]$~--- порог срабатывания;
    \item $\text{act}_i \in \{\texttt{block}, \texttt{transform}, \texttt{pass}\}$~--- политика реагирования при $h_i(x_i) \ge \theta_i$.
\end{itemize}
При $\text{act}_i = \texttt{block}$ возвращается $a_\varnothing$ и каскад завершается. При $\text{act}_i = \texttt{transform}$ вход $x_i$ подменяется на санитизированный $\tilde{x}_i$ и стадия передаёт его дальше.

\subsubsection{Базовые методы защиты: фильтрация и Robust System Prompting}
\label{sec:def_basic}

\paragraph{Фильтрация (Filtering).}
Фильтры~--- вырожденный случай стадий каскада с $\text{act}_i = \texttt{block}$ и сигналом-индикатором.
\begin{itemize}
    \item \textbf{Input Filter:} $F_{in}: \mathcal{V}^* \to \{0, 1\}$. Пропускает запрос $q$, только если $F_{in}(q) = 0$.
    \item \textbf{Data Filter:} $F_{data}: \mathcal{P}(\mathcal{V}^*) \to \mathcal{P}(\mathcal{V}^*)$~--- функция очистки контекста:
    \begin{equation}
        F_{data}(C) = \{ c \in C \mid \text{MaliciousScore}(c) < \epsilon \}
    \end{equation}
\end{itemize}

\paragraph{Robust System Prompting.}
Модификация функции формирования промпта $\mathcal{T}$: вместо стандартного шаблона используется $\mathcal{T}_{robust}$ с защитным системным префиксом $P_{sys}$:
\begin{equation}
    \mathcal{T}_{robust}(q, C) = P_{sys} \oplus \text{Context: } C \oplus \text{Query: } q.
\end{equation}
Цель~--- сместить распределение генерации для вредоносных инструкций так, чтобы условная энтропия генерации на вредоносном входе была высока: $H\bigl(G(\mathcal{T}_{robust}(q_{malicious}, C))\bigr) \to \infty$ (на практике~--- отказ или бессодержательный ответ).

\subsubsection{Формализация SOTA-защит как компонентов каскада}
\label{sec:def_sota}

В данном подразделе описаны три класса современных защитных механизмов, которые в~\S\ref{sec:def_hard} используются в качестве компонентов предлагаемой архитектуры.

\paragraph{1.~RAGFort: полная black-box-адаптация \cite{RAGFort2025}.}

\textit{Механизм.} RAGFort \cite{RAGFort2025} реализует двухкомпонентную защиту: \textbf{Constrained Cascade Generation} (CCG) на уровне ответа и \textbf{Contrastive Re-indexing} (CRI) на уровне корпуса. В настоящей работе оба компонента адаптируются к сценарию black-box доступа к LLM и к работе с уже построенным внешним эмбеддером~--- эта адаптация называется \emph{полной RAGFort-адаптацией} и используется в эксперименте как профиль \texttt{ragfort} (\S\ref{sec:experiment_design}, табл.~\ref{tab:defense_profiles}). Адаптация состоит из двух подкомпонентов:
\begin{itemize}[label=$\circ$]
    \item \textbf{Draft-then-Verify} (DtV)~--- практический суррогат CCG: генератор $G$ формирует черновой ответ $\hat{a}$, после чего отдельная модель-верификатор $V$ принимает или отклоняет его, что заменяет токен-уровневое стробирование оригинального CCG в условиях недоступности логитов LLM. Формально оператор DtV определяется ниже формулами~\eqref{eq:draft}--\eqref{eq:verify}.
    \item \textbf{Topic-Consistent Re-ranking} (TCR)~--- онлайн-аналог CRI: вместо обучения нового контрастивного эмбеддера выполняется оффлайн-кластеризация корпуса в пространстве уже имеющегося эмбеддера и онлайн-переранжирование кандидатов с учётом их соответствия родному топику запроса (формула~\eqref{eq:tcr_score}); тем самым эффект «семантической изоляции тем» из CRI достигается без переобучения. Подробное описание TCR приведено в~\S\ref{sec:def_hard}, поскольку TCR используется как разделяемый компонент: входит и в профиль \texttt{ragfort}, и в профиль \texttt{hard}.
\end{itemize}

\textit{Верификатор для DtV.} Вводится модель-верификатор $V: \mathcal{V}^* \times \mathcal{V}^* \times \mathcal{V}^* \to [0, 1]$, оценивающая условную вероятность отсутствия утечки:
\begin{equation}
    V(q, C, \hat{a}) \approx \mathbb{P}\bigl(\mathbb{I}_{leak}(\hat{a}, S_{priv}) = 0 \mid q, C, \hat{a}\bigr).
\end{equation}
Процесс генерации:
\begin{align}
    \hat{a}\ &=\ G(\mathcal{T}(q, C)) \quad (\text{Draft generation}) \label{eq:draft}\\
    a_{final}\ &=\
    \begin{cases}
      \hat{a}, & V(q, C, \hat{a}) \ge \theta_{safe}, \\
      a_\varnothing, & \text{иначе.}
    \end{cases}
    \label{eq:verify}
\end{align}

\textit{Обоснование эффективности.} Обозначим $P(\mathrm{Leak}_G)$~--- вероятность того, что генератор $G$ создаст утечку, а $P(\mathrm{Miss}_V \mid \mathrm{Leak}_G)$~--- вероятность того, что верификатор $V$ пропустит её (False Negative). Тогда итоговая вероятность успешной атаки для каскада:
\begin{equation}
    P(\text{Success}) = P(\mathrm{Leak}_G) \cdot P(\mathrm{Miss}_V \mid \mathrm{Leak}_G).
\end{equation}
Поскольку верификатор $V$ решает задачу бинарной классификации (утечка / не утечка), вычислительно более простую, чем задача генерации, можно ожидать $P(\mathrm{Miss}_V) \ll 1$, что существенно снижает ASR системы по сравнению с одиночной моделью: $P(\text{Success}) \ll P(\mathrm{Leak}_G)$.

\begin{remark}
В оригинальной работе \cite{RAGFort2025} Constrained Cascade Generation действует на уровне ответа и опирается на \textit{внутреннюю оценку неопределённости} генератора (токен-уровневое стробирование по энтропии распределения следующего токена), а не на отдельную модель-верификатор. В настоящей работе реализуется практический суррогат~--- \textit{Draft-then-Verify} с двумя моделями, что позволяет использовать RAGFort-подход в сценарии black-box API доступа к LLM, когда внутренние состояния генератора недоступны.
\end{remark}

\paragraph{2.~LeakSealer (Behavioral Analysis \cite{LeakSealer2025}).}

\textit{Механизм.} Метод строит «карты использования» (Usage Maps)~--- кластеризацию эмбеддингов легитимных запросов. Пусть $H = \{\mu_1, \dots, \mu_K\}$~--- множество центроидов. Детектор аномалий:
\begin{equation}
    \mathrm{IsAttack}_{\text{LS}}(q) = \mathbb{I}\!\left( \min_{\mu \in H} \lVert E(q) - \mu \rVert_2 > \sigma_{threshold} \right).
\end{equation}

\textit{Обоснование эффективности.} Метод опирается на \textbf{гипотезу многообразия} (Manifold Hypothesis): легитимные запросы $q_{legit}$ сосредоточены на низкоразмерном многообразии $\mathcal{M} \subset \mathbb{R}^k$ в пространстве эмбеддингов. Атакующие запросы, особенно полученные оптимизационными методами (в частности, $O^*_{jail}$ в SECRET), являются статистическими выбросами:
\begin{equation}
    \text{dist}(E(q_{atk}), \mathcal{M}) \gg \text{dist}(E(q_{legit}), \mathcal{M}).
\end{equation}

\begin{remark}[Ограниченная применимость к атакам SECRET]
\label{rem:leaksealer_limit}
В рамках принятой модели угроз LeakSealer формально применим (использует только эмбеддинг запроса). Однако в SOTA-атаке SECRET (\S\ref{sec:secret_formalization}) триггер $T_{retr}$ \textit{целенаправленно оптимизируется} на минимизацию расстояния до центроида целевого тематического кластера (формула~\eqref{eq:cft}), а jailbreak-оператор $O_{jail}$ маскируется под легитимные запросы. Тем самым полный атакующий запрос $x = I_{ext} \oplus O_{jail} \oplus T_{retr}$ конструктивно располагается \textit{внутри} многообразия $\mathcal{M}$, а не в его дополнении. Следовательно, single-query OOD-сигнал LeakSealer теоретически слаб против именно той атаки, для нейтрализации которой строится метод HARD. По этой причине в настоящей работе LeakSealer не реализуется как самостоятельная стадия HARD и не сравнивается численно в главе~3. Идея кластерного анализа эмбеддингов, лежащая в основе LeakSealer, заимствуется в компоненте IRD (\S\ref{sec:def_hard}), но применяется не к запросу целиком, а к ретривальным отпечаткам его семантических фрагментов, что и обеспечивает устойчивость нового сигнала к структуре SECRET-атаки (Предложение~\ref{prop:ird_atk}).
\end{remark}

\paragraph{3.~ControlNET (Activation Monitoring \cite{ControlNET2025}).}

\textit{Механизм.} Анализирует внутренние активации нейронов LLM. Пусть $L_i(p) \in \mathbb{R}^{d_i}$~--- вектор активаций $i$-го слоя при подаче промпта $p$. Защита использует классификатор $\mathrm{Cls}$, действующий на конкатенации активаций выбранных слоёв $k_1, \dots, k_m$:
\begin{equation}
    \text{Block} \iff \mathrm{Cls}\bigl(L_{k_1}(p) \oplus \dots \oplus L_{k_m}(p)\bigr) = \texttt{Malicious}.
\end{equation}

\textit{Обоснование эффективности.} Подход базируется на свойстве \textbf{линейной разделимости семантических концептов} в глубоких слоях трансформеров. Даже если атакующий запрос $q_{atk}$ лексически замаскирован под легитимный, его намерение (intent) проявляется как специфический паттерн активации (\textit{activation shift}), необходимый для переключения модели в режим выдачи запрещённой информации. Предполагается существование разделяющей гиперплоскости $w$:
\begin{equation}
    \exists\, w: \quad \forall\, x \in \mathcal{Q}_{atk},\ w^{\top} \bar{L}(x) > 0,\ \
    \forall\, x \in \mathcal{Q}_{legit},\ w^{\top} \bar{L}(x) < 0,
\end{equation}
где $\bar{L}(x) = L_{k_1}(x) \oplus \dots \oplus L_{k_m}(x)$. Обучение $\mathrm{Cls}$ сводится к поиску такой гиперплоскости.

\begin{remark}[Неприменимость в принятой модели угроз]
\label{rem:controlnet_limit}
Эффективность ControlNET принципиально опирается на доступ к активациям $L_{k_i}(p)$ внутренних слоёв LLM-генератора. В принятой в настоящей работе модели угроз (\S\ref{sec:threat_model}, black-box API доступ к LLM) такой доступ отсутствует, что делает прямое применение метода невозможным. Эмбеддинговый суррогат (детектор смещения в пространстве внешнего эмбеддера $E$ относительно центроида бенигн-распределения) при этом \textit{не используется}: его сигнал по форме совпадает с LeakSealer-style OOD-детектором (см.~\S\ref{rem:leaksealer_limit}), не несёт независимой информации и при включении в каскад приводил бы к двойному учёту одного и того же признака, нарушая допущение условной независимости сигналов в теореме~\ref{thm:cascade_bound}. Поэтому ControlNET в эксперименте не реализуется и в HARD не используется ни в каком виде.
\end{remark}

\subsubsection{Предлагаемый метод: каскадная адаптивная защита с Intent--Retrieval Dissociation (HARD)}
\label{sec:def_hard}

Основной вклад настоящей работы~--- композиционный защитный механизм $\mathrm{Def}^{\text{HARD}}$ (\textit{Hybrid Adaptive RAG Defense}), сочетающий полную RAGFort-адаптацию (DtV + TCR, \S\ref{sec:def_sota}) с тремя оригинальными компонентами, разработанными в настоящей работе:
\begin{itemize}[label=$\bullet$]
    \item \textbf{IRD} (Intent--Retrieval Dissociation)~--- сигнал стадии~1, эксплуатирующий структурную декомпозируемость SOTA-атак извлечения (формула~\eqref{eq:secret_decomp});
    \item \textbf{Output Leak Scanner}~--- детерминированный пост-фильтр стадии~4 (формула~\eqref{eq:leak_scan}), дополняющий LLM-верификатор регулярными выражениями и LCSR-детектором дословного цитирования контекста;
    \item \textbf{Risk-Budget Thresholding}~--- адаптивное порождение порогов стадий каскада на основе агрегированного риска предыдущих стадий (формула~\eqref{eq:risk_budget}).
\end{itemize}
Таким образом, разница между профилями \texttt{ragfort} и \texttt{hard} сводится к перечисленным трём компонентам. TCR в HARD не переименовывается и не модифицируется: используется тот же модуль, что и в профиле \texttt{ragfort}, что обеспечивает корректность сравнения и исключает приписывание метода HARD эффекта от переранжирования.

\paragraph{Архитектура.}
Защитный оператор $\mathrm{Def}^{\text{HARD}}$ представляется как композиция четырёх стадий~\eqref{eq:def_cascade}:
\begin{equation}
    \mathrm{Def}^{\text{HARD}} \;=\;
    \mathrm{Def}_{out}^{(4)} \circ
    \mathrm{Def}_{gen}^{(3)} \circ
    \mathrm{Def}_{ret}^{(2)} \circ
    \mathrm{Def}_{in}^{(1)}.
\end{equation}
Каждая стадия характеризуется своей сигнальной функцией $h_i$, порогом $\theta_i$ и политикой срабатывания (табл.~\ref{tab:hard_stages}).

\begin{table}[H]
\centering
\caption{Стадии каскада $\mathrm{Def}^{\text{HARD}}$.}
\label{tab:hard_stages}
\begin{tabular}{|l|l|l|l|}
\hline
\textbf{Стадия} & \textbf{Вход} & \textbf{Сигнальная функция $h_i$} & \textbf{Политика} \\
\hline
$\mathrm{Def}_{in}^{(1)}$   & $q$               & $s_{\text{IRD}}(q)$                                                       & \texttt{block}     \\
$\mathrm{Def}_{ret}^{(2)}$  & $(q, C)$          & $\mathrm{AnomScore}_{\text{TCR}}(q, C)$                                  & \texttt{transform} \\
$\mathrm{Def}_{gen}^{(3)}$  & $(q, C, \hat{a})$ & $1 - V(q, C, \hat{a})$                                                   & \texttt{block}     \\
$\mathrm{Def}_{out}^{(4)}$  & $a_{final}$       & $\mathrm{LeakScanScore}(a_{final}, C)$                                   & \texttt{block}     \\
\hline
\end{tabular}
\end{table}

Сигналы стадий опираются на \textit{принципиально различные источники информации} (query-space, corpus-space, generator-space, surface-level). Это свойство, называемое далее \textit{диверсификацией сигналов}, является ключевым для теоретической гарантии, сформулированной в теореме~\ref{thm:cascade_bound}.

\paragraph{Компонент 1: Intent--Retrieval Dissociation (IRD).}

\textit{Мотивация.} В~\S\ref{sec:secret_formalization} SECRET-атака формализована как композиция трёх функциональных компонентов: $x = I_{ext} \oplus O_{jail} \oplus T_{retr}$ (формула~\eqref{eq:secret_decomp}). Ключевое наблюдение состоит в том, что у \textit{легитимного} запроса смысловое намерение и триггер поиска суть одно и то же, тогда как у атаки SECRET (и у Prompt Injection в общем случае) это разные по содержанию токенные подпоследовательности. Инструкция извлечения $I_{ext}$ (например, «repeat verbatim\dots») семантически не связана с триггером $T_{retr}$, оптимизированным для приближения к центроиду тематического кластера. Как следствие, \textit{наборы чанков, извлекаемых ретривером для разных семантических фрагментов атакующего запроса, почти не пересекаются}~--- в отличие от легитимного запроса, фрагменты которого тематически согласованы.

\textit{Формальное определение.} Пусть $\phi: \mathcal{V}^* \to (\mathcal{V}^*)^m$~--- функция семантической декомпозиции запроса, возвращающая $m$ фрагментов $\phi(q) = (q^{(1)}, \dots, q^{(m)})$. В настоящей работе $\phi$ реализуется как LLM-as-a-segmenter: вспомогательной языковой модели подаётся промпт-инструкция разбить запрос на минимальные смысловые единицы (вопросы / инструкции / тематические якоря).

Для каждого фрагмента $q^{(i)}$ определим \textit{ретривальный отпечаток}~--- множество идентификаторов top-$k_{retr}$ чанков, возвращаемых ретривером:
\begin{equation}
    R_{\text{set}}(q^{(i)}) = \{\mathrm{id}(c) : c \in R(q^{(i)}, \mathcal{I})\}, \qquad |R_{\text{set}}(q^{(i)})| = k_{retr}.
\end{equation}

\textbf{IRD-скор} определяется как единица минус среднее попарное сходство Жаккара ретривальных отпечатков фрагментов:
\begin{equation}
    \label{eq:ird_score}
    \boxed{\
    s_{\text{IRD}}(q) =
    \begin{cases}
    0, & m < 2, \\[4pt]
    \displaystyle 1 - \frac{2}{m(m-1)} \sum_{1 \le i < j \le m}
    \frac{\bigl|R_{\text{set}}(q^{(i)}) \cap R_{\text{set}}(q^{(j)})\bigr|}{\bigl|R_{\text{set}}(q^{(i)}) \cup R_{\text{set}}(q^{(j)})\bigr|},
    & m \ge 2.
    \end{cases}
    \ }
\end{equation}

Значения $s_{\text{IRD}} \in [0, 1]$: нулю соответствуют полностью согласованные фрагменты (все тянут одни и те же чанки), единице~--- полная дезинтеграция (фрагменты тянут попарно непересекающиеся подмножества корпуса).

\begin{proposition}[Свойство IRD для SECRET-атак]
\label{prop:ird_atk}
Пусть для атакующего запроса $x = I_{ext} \oplus O_{jail} \oplus T_{retr}$ декомпозиция $\phi$ разделяет три функциональных компонента: $\phi(x) = (I_{ext}, O_{jail}, T_{retr})$. Тогда при независимо сгенерированных $I_{ext}$ и $T_{retr}$ и при равномерной кластеризации корпуса на $K$ кластеров выполняется
\begin{equation}
    \mathbb{E}[s_{\text{IRD}}(x)] \ge 1 - \frac{1}{K}.
\end{equation}
\end{proposition}
\begin{proof}[Эскиз доказательства]
$I_{ext}$ является инструкцией, тематически не связанной с корпусом, тогда как $T_{retr}$ оптимизирован для целевого кластера $\mu_i$. При равномерной кластеризации вероятность того, что top-$k_{retr}$ по $I_{ext}$ и top-$k_{retr}$ по $T_{retr}$ попадут в один кластер, не превосходит $1/K$. Отсюда $\mathbb{E}[\mathrm{Jaccard}] \le 1/K$ и $\mathbb{E}[s_{\text{IRD}}] \ge 1 - 1/K$.
\end{proof}

\begin{proposition}[Свойство IRD для легитимных запросов]
\label{prop:ird_legit}
Пусть $q \in \mathcal{Q}_{legit}$~--- тематически однородный запрос, и $\phi$ производит декомпозицию на грамматически разные, но тематически связные фрагменты. Тогда при достаточном $k_{retr}$ существует $\eta > 0$ такое, что
\begin{equation}
    \mathbb{E}[s_{\text{IRD}}(q)] \le 1 - \eta.
\end{equation}
\end{proposition}
Эмпирическая проверка Предложений~\ref{prop:ird_atk} и \ref{prop:ird_legit} выполняется в главе~3 по метрике $\mathrm{AUC}_{\text{IRD}}$ (формула~\eqref{eq:auc_ird}).

\textit{Решающее правило.} Запрос помечается как подозрительный стадией $\mathrm{Def}_{in}^{(1)}$, если выполняется
\begin{equation}
    \label{eq:stage1_rule}
    h_1(q) = s_{\text{IRD}}(q) \ge \theta_1.
\end{equation}
Порог $\theta_1$ калибруется по бенигн-выборке как квантиль распределения $s_{\text{IRD}}$ на легитимных запросах (по умолчанию~--- $0.95$-квантиль). Подключение дополнительного OOD-сигнала $s_{\text{OOD}}$ в духе LeakSealer на этой стадии \textit{не выполняется} по причинам, указанным в~\S\ref{rem:leaksealer_limit}: такой сигнал теоретически слаб против атак SECRET и нарушает диверсификацию сигналов стадий, на которой основана граница из теоремы~\ref{thm:cascade_bound}.

\paragraph{Компонент 2: Topic-Consistent Re-ranking (TCR).}

Онлайн-версия идеи Contrastive Re-indexing из RAGFort \cite{RAGFort2025}, адаптированная под чёрный ящик ретривера. Компонент является \textit{разделяемым} между профилями \texttt{ragfort} и \texttt{hard} (см.~\S\ref{sec:def_sota}, \S\ref{sec:experiment_design}): он не относится к оригинальному вкладу настоящей работы, а реализует black-box-адаптацию RAGFort.

\textit{Оффлайн-этап.} На этапе построения индекса корпус $\mathcal{D}$ разбивается на $K_{top}$ тематических кластеров методом $k$-means в пространстве эмбеддингов; хранятся центроиды $\{\mu_1, \dots, \mu_{K_{top}}\}$ и метка топика $\mathrm{topic}(c) \in \{1, \dots, K_{top}\}$ для каждого чанка.

\textit{Онлайн-этап.} Для запроса $q$ определяется родной топик $k^*(q) = \arg\max_k \mathrm{sim}(E(q), \mu_k)$, после чего каждый кандидат $c \in R(q, \mathcal{I})$ переранжируется по скорректированной оценке:
\begin{equation}
    \label{eq:tcr_score}
    \mathrm{score}(c \mid q) = \alpha\,\mathrm{sim}(E(q), E(c)) + \beta\,\mathrm{sim}(E(c), \mu_{k^*(q)}) - \gamma\,\mathrm{anom}(c),
\end{equation}
где $\alpha + \beta = 1$, $\gamma \ge 0$, а
\begin{equation}
    \mathrm{anom}(c) = \bigl\lVert E(c) - \mu_{\mathrm{topic}(c)} \bigr\rVert_2
\end{equation}
---~\textit{собственная аномальность чанка} относительно его родного кластера. Компонент $\mathrm{anom}(c)$ подавляет poison-чанки (они часто имеют большое расстояние до кластера-основы либо искусственно сдвинуты к чужому триггер-кластеру).

Сигнал стадии:
\begin{equation}
    h_2(q, C) = \frac{1}{|C|} \sum_{c \in C} \mathrm{anom}(c).
\end{equation}
Политика: переранжирование по~\eqref{eq:tcr_score} и отбрасывание чанков с $\mathrm{anom}(c) > \theta_2$.

\paragraph{Компонент 3: Draft-then-Verify (каскадная генерация).}

Совпадает с \eqref{eq:draft}--\eqref{eq:verify} (см.~\S\ref{sec:def_sota}); как и TCR, является разделяемым компонентом полной RAGFort-адаптации (профиль \texttt{ragfort}) и метода HARD. Стадия имеет сигнал
\begin{equation}
    h_3(q, C, \hat{a}) = 1 - V(q, C, \hat{a}).
\end{equation}

\paragraph{Компонент 4: Output Leak Scanner.}

Стадия $\mathrm{Def}_{out}^{(4)}$~--- детерминированный пост-фильтр, дополняющий LLM-верификатор. Вводится скан-функция
\begin{equation}
    \label{eq:leak_scan}
    h_4(a_{final}, C, \mathcal{P}) = \max\!\left(
      \mathrm{RegexHit}(a_{final}, \mathcal{P}),\
      \max_{c \in C} \mathbb{I}\!\left[\mathrm{LCSR}(a_{final}, c) > \tau_{copy}\right]
    \right),
\end{equation}
где $\mathcal{P}$~--- множество предопределённых регулярных выражений чувствительных данных (ключи, пароли, идентификаторы, триггер-строки backdoor-атаки), $\mathrm{LCSR}$~--- нормированная длина наибольшей общей подпоследовательности ответа и чанка (защита от дословного цитирования контекста), $\tau_{copy}$~--- порог цитирования.

\paragraph{Адаптивные пороги: Risk-Budget Propagation.}

Наивная версия каскада использует статические пороги $\theta_i^{(0)}$. Предлагаемая модификация~--- \textit{распространение риск-бюджета}: порог каждой последующей стадии понижается пропорционально агрегированному риску, накопленному на предыдущих стадиях:
\begin{equation}
    \label{eq:risk_budget}
    \theta_i(q) = \theta_i^{(0)} - \lambda_i \sum_{j < i} r_j(q),
\end{equation}
где $r_j(q) \in [0, 1]$~--- нормированный остаточный риск стадии $j$ (например, $r_j = h_j(\cdot) / \theta_j^{(0)}$ при $h_j < \theta_j^{(0)}$), а $\lambda_i \ge 0$~--- настраиваемые коэффициенты связанности.

\textit{Интерпретация.} Если стадия $j$ «почти сработала», но не преодолела порог, то стадия $i > j$ становится строже. Это смягчает проблему жёстких порогов: атакующий, научившийся оставаться чуть ниже каждого порога по отдельности, сталкивается с суммарным увеличением чувствительности каскада.

\paragraph{Теоретическое обоснование: граница ASR.}

\begin{theorem}[Cascade Bound]
\label{thm:cascade_bound}
Пусть для механизма $\mathrm{Def}^{\text{HARD}}$ на тестовом множестве атак $\mathcal{Q}_{atk}$ определены индивидуальные вероятности пропуска стадий:
\begin{equation}
    m_i(\theta_i) = \mathbb{P}\bigl[ h_i < \theta_i \,\bigm|\, q \in \mathcal{Q}_{atk},\ \text{атака достигла стадии } i \bigr], \quad i = 1, 2, 3, 4.
\end{equation}
Если сигналы стадий условно независимы относительно события утечки, то
\begin{equation}
    \label{eq:cascade_bound}
    \mathrm{ASR}\bigl(\mathrm{Def}^{\text{HARD}}\bigr) \le \prod_{i=1}^{4} m_i(\theta_i).
\end{equation}
\end{theorem}
\begin{proof}
Обозначим $A$ событие успешной атаки (фиксированный $q \in \mathcal{Q}_{atk}$ приводит к ответу $a$ с $\mathbb{I}_{leak}(a, S_{priv}) = 1$). Для утечки необходимо, чтобы все четыре стадии «пропустили» атаку; обозначим $M_i = \{h_i < \theta_i\}$. Тогда
\begin{equation}
    \mathbb{P}(A) \le \mathbb{P}(M_1 \cap M_2 \cap M_3 \cap M_4).
\end{equation}
По определению условной вероятности и в предположении условной независимости сигналов относительно $q$:
\begin{equation}
    \mathbb{P}(M_1 \cap M_2 \cap M_3 \cap M_4) = \prod_{i=1}^4 \mathbb{P}(M_i) = \prod_{i=1}^4 m_i(\theta_i),
\end{equation}
откуда и следует~\eqref{eq:cascade_bound}.
\end{proof}

\begin{corollary}[Строгое убывание ASR в адаптивном режиме]
\label{cor:adaptive}
Пусть функции $m_i(\theta)$ монотонно возрастают по $\theta$ (чем выше порог, тем выше вероятность пропуска), и пусть на типовой атаке $q$ все $r_j(q) > 0$. Тогда для адаптивных порогов~\eqref{eq:risk_budget} с $\lambda_i > 0$ справедливо
\begin{equation}
    \prod_{i=1}^{4} m_i\!\bigl(\theta_i(q)\bigr) < \prod_{i=1}^{4} m_i\!\bigl(\theta_i^{(0)}\bigr),
\end{equation}
т.\,е. распространение риск-бюджета строго улучшает верхнюю границу ASR относительно статических порогов при условии ненулевых остаточных рисков.
\end{corollary}

\begin{remark}
Условная независимость сигналов стадий в общем случае не выполняется строго: например, сильный OOD-сигнал коррелирует с «подозрительным» черновиком. Теорема~\ref{thm:cascade_bound} поэтому имеет характер \textit{архитектурного принципа проектирования}: каскад должен состоять из стадий, опирающихся на максимально разнородные источники сигнала (query-space, corpus-space, generator-space, surface-level). Эмпирическая проверка степени зависимости между стадиями проводится в экспериментальной главе по метрикам $\mathrm{StageCorrelation}$ и $\kappa$ (формулы~\eqref{eq:stage_corr}, \eqref{eq:kappa}).
\end{remark}

\subsection{Метрики оценки}
\label{sec:metrics}

Для количественной оценки эффективности атак и защитных механизмов вводятся следующие метрики. Первые четыре~--- стандартные метрики для сравнения произвольных защит; последние четыре~--- вспомогательные, необходимые для эмпирической проверки допущений каскадной модели HARD.

\begin{enumerate}
    \item \textbf{Attack Success Rate (ASR):} доля успешных атак на тестовом наборе $\mathcal{Q}_{atk}$:
    \begin{equation}
        \label{eq:asr}
        \mathrm{ASR} = \frac{1}{|\mathcal{Q}_{atk}|} \sum_{q \in \mathcal{Q}_{atk}} \mathbb{I}_{leak}\bigl(S(q), S_{priv}\bigr).
    \end{equation}

    \item \textbf{False Positive Rate (FPR):} доля легитимных запросов из $\mathcal{Q}_{legit}$, ошибочно заблокированных защитой:
    \begin{equation}
        \label{eq:fpr}
        \mathrm{FPR} = \frac{1}{|\mathcal{Q}_{legit}|} \sum_{q \in \mathcal{Q}_{legit}} \mathbb{I}\bigl(\mathrm{Def}(q) = \texttt{Block}\bigr).
    \end{equation}

    \item \textbf{Benign Performance Drop (BPD):} снижение качества работы системы на легитимных запросах при включении защиты:
    \begin{equation}
        \label{eq:bpd}
        \mathrm{BPD} = \mathrm{Score}_{benign}(\mathcal{S}_{orig}) - \mathrm{Score}_{benign}(\mathcal{S}_{safe}),
    \end{equation}
    где $\mathrm{Score}$ может быть метрикой ROUGE, BERTScore или точностью QA.

    \item \textbf{Latency Overhead:} накладные расходы по времени выполнения:
    \begin{equation}
        \label{eq:latency}
        \Delta T = \frac{1}{|\mathcal{Q}|} \sum_{q \in \mathcal{Q}} \bigl(T_{safe}(q) - T_{orig}(q)\bigr).
    \end{equation}

    \item \textbf{Stage Pass Rate.} Для каждой стадии $i \in \{1, 2, 3, 4\}$ каскада $\mathrm{Def}^{\text{HARD}}$ определим долю атакующих запросов, прошедших стадию без блокировки или трансформации:
    \begin{equation}
        \label{eq:stage_pass}
        \mathrm{StagePassRate}_i =
        \frac{1}{|\mathcal{Q}_{atk}|}
        \sum_{q \in \mathcal{Q}_{atk}}
        \mathbb{I}\bigl[h_i(x_i(q)) < \theta_i\bigr].
    \end{equation}
    Метрика позволяет идентифицировать «узкие места» каскада и оценить эмпирическую вероятность пропуска $m_i$, фигурирующую в теореме~\ref{thm:cascade_bound}: $\widehat{m}_i = \mathrm{StagePassRate}_i$.

    \item \textbf{Stage Correlation.} Для эмпирической проверки предположения условной независимости (см. замечание к теореме~\ref{thm:cascade_bound}) вводится метрика попарной корреляции бинарных флагов срабатывания стадий, усреднённая по $\mathcal{Q}_{atk}$:
    \begin{equation}
        \label{eq:stage_corr}
        \mathrm{StageCorr}(i, j) = \mathrm{Corr}\!\bigl(\mathbb{I}[h_i \ge \theta_i],\ \mathbb{I}[h_j \ge \theta_j]\bigr), \qquad i \ne j.
    \end{equation}
    Низкие значения $|\mathrm{StageCorr}| \to 0$ подтверждают архитектурный принцип диверсификации сигналов.

    \item \textbf{IRD Separation (ROC-AUC).} Дискриминативная сила предлагаемого IRD-сигнала:
    \begin{equation}
        \label{eq:auc_ird}
        \mathrm{AUC}_{\text{IRD}} = \mathbb{P}\bigl(s_{\text{IRD}}(q_{atk}) > s_{\text{IRD}}(q_{legit})\bigr),
    \end{equation}
    где пара $(q_{atk}, q_{legit})$ берётся случайно из $\mathcal{Q}_{atk} \times \mathcal{Q}_{legit}$. Значение $\mathrm{AUC}_{\text{IRD}} = 0{,}5$ соответствует отсутствию дискриминации, $\mathrm{AUC}_{\text{IRD}} = 1$~--- идеальному разделению.

    \item \textbf{Cascade Bound Tightness.} Эмпирическая плотность теоретической границы~\eqref{eq:cascade_bound}:
    \begin{equation}
        \label{eq:kappa}
        \kappa = \frac{\mathrm{ASR}_{emp}}{\prod_{i=1}^{4} \widehat{m}_i}, \qquad \widehat{m}_i = \mathrm{StagePassRate}_i.
    \end{equation}
    Значение $\kappa \approx 1$ означает, что предположение условной независимости выполняется эмпирически; $\kappa < 1$~--- стадии работают «лучше независимых» (положительная синергия).
\end{enumerate}

\newpage
\section{Требования к программному комплексу и дизайн эксперимента}
\label{sec:software_complex}

Целью разрабатываемого программного комплекса является воспроизводимая оценка качества и защищённости RAG-систем от атак извлечения данных, а также экспериментальная проверка предлагаемого метода HARD. Комплекс должен поддерживать полный цикл исследования: построение RAG-системы на заданном корпусе, измерение качества ответов без атак, моделирование атак извлечения, подключение защитных профилей и повторное измерение качества и устойчивости. В настоящей главе формулируются требования к комплексу, его архитектура и дизайн экспериментальной оценки.

\subsection{Требования к программному комплексу}
\label{sec:software_requirements}

Требования разделены на функциональные (FR), нефункциональные (NFR) и требования к воспроизводимости. Функциональные требования определяют набор сценариев, которые должен выполнять комплекс; нефункциональные фиксируют ограничения вычислительной среды и качества реализации; требования к воспроизводимости необходимы для корректной интерпретации экспериментальных результатов и независимой проверки.

\subsubsection{Функциональные требования}

\begin{table}[H]
\centering
\caption{Функциональные требования к программному комплексу.}
\label{tab:functional_requirements}
\begin{tabular}{|p{0.10\textwidth}|p{0.82\textwidth}|}
\hline
\textbf{ID} & \textbf{Требование} \\
\hline
FR-1 & Сборка RAG-системы. Комплекс должен по заданному корпусу документов $\mathcal{D}$ выполнять разбиение на чанки, вычисление эмбеддингов $E$, построение векторного индекса $\mathcal{I}$ и его сохранение в неизменяемом виде на диск. \\
\hline
FR-2 & Измерение базового качества. Комплекс должен на легитимном наборе вопросов $\mathcal{Q}_{legit}$ без атак и без защит измерять $\mathrm{Score}_{benign}(\mathcal{S}_{orig})$ через метрики lexical overlap, token-F1 и LLM-as-a-judge (\S\ref{sec:metrics}). \\
\hline
FR-3 & Моделирование атак извлечения. Комплекс должен моделировать три класса атак извлечения данных: Prompt Injection (\S\ref{sec:pi_formalization}), Backdoor/Data Poisoning (\S\ref{sec:backdoor_formalization}) и SECRET как SOTA-атаку (\S\ref{sec:secret_formalization}) с явной декомпозицией $x = I_{ext} \oplus O_{jail} \oplus T_{retr}$. \\
\hline
FR-4 & Измерение устойчивости. Комплекс должен на каждом классе атак из FR-3 вычислять Attack Success Rate (формула~\eqref{eq:asr}) с использованием детектора утечки $\mathbb{I}_{leak}$ в режимах substring, regex и LLM-as-a-judge. \\
\hline
FR-5 & Профили защиты. Комплекс должен поддерживать четыре профиля защиты, активируемых по имени: \texttt{none}, \texttt{basic-filters}, \texttt{ragfort}, \texttt{hard}. Профиль \texttt{ragfort} соответствует \emph{полной} black-box-адаптации архитектуры RAGFort~\cite{RAGFort2025}: Constrained Cascade Generation в виде Draft-then-Verify $+$ Contrastive Re-indexing в виде Topic-Consistent Re-ranking (\S\ref{sec:def_sota}). Этот профиль служит главным baseline для оценки вклада оригинальных компонентов HARD (IRD, Output Leak Scanner, Risk-Budget Thresholding) и реализуется на тех же модулях, что и соответствующие подкомпоненты HARD. \\
\hline
FR-6 & Повторное измерение со включёнными защитами. Для каждого профиля защиты $p$ комплекс должен повторно измерять ASR, FPR (формула~\eqref{eq:fpr}), BPD (формула~\eqref{eq:bpd}) и Latency Overhead (формула~\eqref{eq:latency}). \\
\hline
FR-7 & Трассировка каскада. Для профиля \texttt{hard} комплекс должен сохранять для каждого запроса полную трассировку стадий: значения сигналов $h_i$, эффективные пороги $\theta_i(q)$ с учётом Risk-Budget, бинарные флаги прохождения стадий и идентификатор стадии блокировки при её наличии. \\
\hline
FR-8 & Эмпирическая проверка теоретической границы. Комплекс должен вычислять вспомогательные метрики StagePassRate (\eqref{eq:stage_pass}), StageCorrelation (\eqref{eq:stage_corr}), $\mathrm{AUC}_{\text{IRD}}$ (\eqref{eq:auc_ird}) и $\kappa$ (\eqref{eq:kappa}) для эмпирической проверки допущений теоремы~\ref{thm:cascade_bound}. \\
\hline
FR-9 & Отчёты и визуализация. Комплекс должен формировать машинно-читаемые отчёты в форматах JSON и CSV, а также графические артефакты (heatmap ASR, гистограммы FPR/BPD/Latency, ROC-кривая IRD, stage-pass breakdown) для включения в экспериментальную главу диссертации. \\
\hline
FR-10 & Пользовательский интерфейс. Комплекс должен предоставлять графический интерфейс пользователя (Web-UI), позволяющий без правки исходного кода: (а)~выбрать корпус и параметры индексации; (б)~выбрать векторное хранилище из числа поддерживаемых; (в)~назначить LLM-провайдер и модель для каждой роли (\texttt{GENERATOR}, \texttt{VERIFIER}, \texttt{JUDGE}, \texttt{SEGMENTER}); (г)~выбрать профиль защиты и тип атаки; (д)~задать число прогонов и значение seed; (е)~запустить эксперимент и наблюдать прогресс; (ж)~просмотреть итоговую таблицу метрик (ASR, FPR, BPD, $\Delta T$) и скачать артефакты прогона (\texttt{summary.json}, \texttt{provenance.json}). Допустимая реализация~--- Streamlit или Gradio. \\
\hline
\end{tabular}
\end{table}

\subsubsection{Нефункциональные требования}

\begin{table}[H]
\centering
\caption{Нефункциональные требования и ограничения.}
\label{tab:nonfunctional_requirements}
\begin{tabular}{|p{0.10\textwidth}|p{0.82\textwidth}|}
\hline
\textbf{ID} & \textbf{Требование} \\
\hline
NFR-1 & Вычислительный бюджет. Полный экспериментальный прогон должен выполняться на портативной рабочей станции без выделенного GPU (характеристики: $\geq 8$ ядер CPU, $\geq 16$~ГБ ОЗУ, ARM- или x86-архитектура) за время не более 1--2 часов. Целевая платформа разработки~--- Apple~Silicon (M-серия). \\
\hline
NFR-2 & Двухрежимный запуск. Комплекс должен поддерживать два режима исполнения: \texttt{full} для финальных результатов и \texttt{fast} для отладки. Режимы должны отличаться только размерами выборок и числом итераций оптимизации, не нарушая логику эксперимента. \\
\hline
NFR-3 & Кэширование LLM. Все LLM-вызовы должны кэшироваться на диске по ключу $\mathrm{SHA256}(\text{provider}, \text{model}, \text{prompt}, \text{temperature}, \text{max\_tokens}, \text{role})$. Инвалидация кэша происходит при изменении любой компоненты ключа; смена модели или роли LLM автоматически порождает новый ключ. \\
\hline
NFR-4 & Детерминизм некритичных вызовов. Все LLM-вызовы в ролях \texttt{VERIFIER}, \texttt{JUDGE}, \texttt{SEGMENTER} должны выполняться с температурой $T = 0$ или эквивалентной минимальной для повышения воспроизводимости. \\
\hline
NFR-5 & Black-box к LLM. Комплекс не должен требовать доступа к внутренним состояниям LLM-генератора (активациям, логитам, весам, градиентам) в соответствии с моделью угроз~\S\ref{sec:threat_model}. \\
\hline
NFR-6 & Модульность. Атаки, защиты, метрики, RAG-пайплайн и CLI-оркестратор должны быть слабо связаны и заменяемы независимо. Зависимости между подсистемами идут через явные интерфейсы (Python-протоколы), внедрение зависимостей~--- через конструкторы. \\
\hline
NFR-7 & Унифицированный доступ к LLM. Комплекс должен предоставлять единый клиентский интерфейс с поддержкой как минимум двух семейств провайдеров (Ollama и OpenAI-совместимый API: Gemini, OpenAI, Together, Anthropic) и допускать независимую конфигурацию модели и провайдера для каждой роли LLM, чтобы дорогие <<умные>> вызовы (\texttt{VERIFIER}, \texttt{JUDGE}) и массовые <<дешёвые>> вызовы (\texttt{GENERATOR}, \texttt{SEGMENTER}) могли быть назначены разным моделям. \\
\hline
NFR-8 & Контроль выборок. При случайной выборке (например, бенигн-тренинг для калибровки порогов) seed фиксируется и логируется. Все стохастические алгоритмы (k-means, выборка триггеров SECRET) принимают параметр \texttt{random\_state}. \\
\hline
\end{tabular}
\end{table}

\subsubsection{Воспроизводимость}

Каждый прогон сопровождается файлом \texttt{provenance.json}, фиксирующим: идентификатор git-коммита, версию корпуса (хеш по содержимому файлов), параметры индекса (top-$k$, размер чанка, эмбеддер), названия и версии моделей по каждой роли LLM, значения seed, активный профиль защиты, параметры атаки и список переменных окружения без секретных значений. Это позволяет отделить изменение результата, вызванное модификацией метода, от изменения результата, вызванного другой моделью или другим корпусом, а также сделать любой прогон независимо проверяемым.

\subsection{Архитектура программного комплекса}
\label{sec:software_architecture}

\subsubsection{Технологический стек}

Программный комплекс реализуется на языке Python 3.12. Ключевые внешние зависимости приведены в таблице~\ref{tab:tech_stack}.

\begin{table}[H]
\centering
\caption{Основные компоненты технологического стека.}
\label{tab:tech_stack}
\begin{tabular}{|p{0.30\textwidth}|p{0.62\textwidth}|}
\hline
\textbf{Компонент} & \textbf{Реализация} \\
\hline
Эмбеддер & \texttt{sentence-transformers}, модель \texttt{paraphrase-multilingual-MiniLM-L12-v2} (размерность $k = 384$) \\
\hline
Векторное хранилище & \texttt{Qdrant} (production-grade vector DB): основной режим~--- локальный сервер, поднимаемый через \texttt{docker compose}; для unit-тестов~--- embedded-режим \texttt{qdrant-client}. Метрика~--- косинусная. \\
\hline
LLM-клиент & Единый интерфейс \texttt{LLMClient(role, prompt, ...)} с двумя реализациями: \texttt{OllamaLLM} и \texttt{OpenAICompatLLM} (Gemini, OpenAI, Together, Anthropic). Провайдер и модель выбираются независимо для каждой роли через \texttt{.env}/YAML. \\
\hline
Модели по умолчанию & \texttt{GENERATOR}: Llama~3.1~8B Instruct (Ollama, локально) либо Gemini~2.0~Flash; \texttt{VERIFIER}, \texttt{JUDGE}: Gemini~2.0~Flash через OpenAI-compat API (требуется <<умная>> модель для надёжной оценки утечки и качества draft); \texttt{SEGMENTER}: Gemini~2.0~Flash либо Llama~3.1~8B. \\
\hline
Кластеризация (TCR, SECRET) & \texttt{scikit-learn} (\texttt{KMeans}) \\
\hline
RAGAS-метрики & \texttt{ragas} для оценки качества baseline RAG \\
\hline
Кэш LLM & on-disk JSON-кэш в каталоге \texttt{runs/llm\_cache/} \\
\hline
Графика & \texttt{matplotlib} (PNG и PDF, без внешних бэкендов) \\
\hline
CLI и orchestration & \texttt{typer} либо \texttt{argparse}; конфигурация через \texttt{.env} и YAML \\
\hline
\end{tabular}
\end{table}

\subsubsection{Компонентная схема}

Программный комплекс организован в виде четырёхслойной модульной архитектуры (рис.~\ref{fig:software_architecture}): \textit{Data Layer} (управление корпусом и индексами), \textit{Pipeline Layer} (RAG-пайплайн, атаки, защиты, оркестратор), \textit{Reporting Layer} (метрики и визуализация) и \textit{Presentation Layer} (Web-UI и CLI). Presentation Layer является \emph{тонкой надстройкой} над Pipeline и Reporting Layers и не дублирует их логику: UI и CLI вызывают одни и те же функции \texttt{Experiment Orchestrator}, обеспечивая идентичность результатов между запуском из интерфейса и из командной строки. Внешние LLM-сервисы, к которым комплекс обращается только через прикладной интерфейс, расположены за границей системы и помечены как black-box.

\begin{figure}[H]
\centering
\begin{tikzpicture}[
    node distance=8mm and 6mm,
    every node/.style={font=\small},
    block/.style={rectangle, draw, rounded corners=2pt, minimum height=8mm, minimum width=32mm, align=center, fill=white},
    bigblock/.style={rectangle, draw, rounded corners=2pt, minimum height=8mm, minimum width=46mm, align=center, fill=white},
    layer/.style={rectangle, draw, dashed, rounded corners=4pt, inner sep=4mm, fill=gray!5},
    arrow/.style={-{Latex[length=2mm]}, thick},
    extlabel/.style={font=\footnotesize\itshape}
]

\node[block] (corpus) {Corpus\\Manager};
\node[block, right=of corpus] (index) {Index\\Builder};
\node[block, right=of index] (poison) {Poisoned-Index\\Builder};
\begin{scope}[on background layer]
\node[layer, fit=(corpus)(index)(poison), label={[anchor=north west]north west:\textbf{Data Layer}}] (datalayer) {};
\end{scope}

\node[bigblock, below=18mm of index] (rag) {RAG Pipeline\\\scriptsize $q \to R(q,\mathcal{I}) \to G(\mathcal{T}(q,C))$};
\node[block, left=of rag] (attacks) {Attack Suite\\\scriptsize PI / Backdoor / SECRET};
\node[block, right=of rag] (defense) {Defense Manager\\\scriptsize none / filters / ragfort / hard};
\node[block, below=8mm of defense] (cascade) {HARD Cascade\\\scriptsize IRD, TCR, DtV, Scanner};
\node[block, below=8mm of attacks] (orch) {Experiment\\Orchestrator};
\begin{scope}[on background layer]
\node[layer, fit=(attacks)(rag)(defense)(cascade)(orch), label={[anchor=north west]north west:\textbf{Pipeline Layer}}] (pipelayer) {};
\end{scope}

\node[block, below=18mm of orch] (metrics) {Metrics Engine\\\scriptsize ASR, FPR, BPD, $\Delta T$, \dots};
\node[block, right=of metrics] (reports) {Report Builder\\\scriptsize summary.json/csv};
\node[block, right=of reports] (plots) {Plot Generator\\\scriptsize PNG/PDF};
\begin{scope}[on background layer]
\node[layer, fit=(metrics)(reports)(plots), label={[anchor=north west]north west:\textbf{Reporting Layer}}] (replayer) {};
\end{scope}

\node[bigblock, fill=gray!15, right=14mm of defense, yshift=-12mm, draw=gray!60] (llm) {External LLM API\\\scriptsize \textit{black-box}};

\draw[arrow] (corpus) -- (index);
\draw[arrow] (corpus) -- (poison);
\draw[arrow] (index) -- (rag);
\draw[arrow] (poison) to[bend right=15] (rag.north);
\draw[arrow] (attacks) -- (rag);
\draw[arrow] (defense) -- (rag);
\draw[arrow] (cascade) -- (defense);
\draw[arrow] (orch.north west) to[bend left=10] (attacks.south);
\draw[arrow] (orch.north) -- (rag.south);
\draw[arrow] (orch.north east) to[bend right=10] (defense.south);
\draw[arrow] (rag.east) to[bend left=20] (llm.west);
\draw[arrow] (cascade.east) to[bend right=10] (llm.south west);
\draw[arrow] (orch) -- (metrics);
\draw[arrow] (metrics) -- (reports);
\draw[arrow] (reports) -- (plots);

\end{tikzpicture}
\caption{Компонентная схема программного комплекса. Сплошные стрелки~--- основные потоки данных и управления. Внешний LLM-сервис рассматривается как black-box в соответствии с~\S\ref{sec:threat_model}.}
\label{fig:software_architecture}
\end{figure}

\subsubsection{Модули}

\begin{table}[H]
\centering
\caption{Модули программного комплекса и их назначение.}
\label{tab:software_modules}
\begin{tabular}{|p{0.24\textwidth}|p{0.66\textwidth}|}
\hline
\textbf{Модуль} & \textbf{Назначение} \\
\hline
\texttt{Corpus Manager} & Чтение и нормализация документов корпуса, разбиение на чанки, объединение базового корпуса и внешних источников, формирование poisoned-корпуса для Backdoor-эксперимента. \\
\hline
\texttt{Index Builder} & Вычисление эмбеддингов чанков, построение и сохранение чистого индекса $\mathcal{I}$ с метаданными чанков. \\
\hline
\texttt{Poisoned-Index Builder} & Сборка отдельного индекса $\mathcal{I}_{pois}$ на корпусе, включающем отравленный документ; используется только для Backdoor-сценария. \\
\hline
\texttt{RAG Pipeline} & Реализация цикла $q \to R(q,\mathcal{I}) \to G(\mathcal{T}(q,C))$ с настраиваемым top-$k$ и системным промптом. \\
\hline
\texttt{Attack Suite} & Генерация атакующих запросов: PI (шаблоны и комбинаторика), Backdoor (триггер $\tau$ + payload $\pi$), SECRET ($I_{ext}$, LLM-as-Optimizer для $O_{jail}$, Cluster-Focused Triggering для $T_{retr}$). \\
\hline
\texttt{Defense Manager} & Активация одного из четырёх профилей защиты по имени; единый интерфейс \texttt{handle(query)} \(\to\) \texttt{ResponseTrace}. \\
\hline
\texttt{HARD Cascade} & Реализация четырёх стадий метода HARD (IRD, TCR, Draft-then-Verify, Output Leak Scanner) и Risk-Budget Thresholding (формула~\eqref{eq:risk_budget}). \\
\hline
\texttt{Metrics Engine} & Вычисление ASR, FPR, BPD, Latency Overhead, StagePassRate, StageCorrelation, $\mathrm{AUC}_{\text{IRD}}$, $\kappa$. \\
\hline
\texttt{Experiment Orchestrator} & CLI-оркестратор: перебор пар (профиль защиты, атака), управление seed-ами, кэшем LLM, путями к индексам и сохранением артефактов. \\
\hline
\texttt{Report Builder} & Сборка \texttt{summary.json}, \texttt{summary.csv}, \texttt{report.json} по каждому прогону, \texttt{experiment.log} и \texttt{provenance.json}. \\
\hline
\texttt{Plot Generator} & Построение графиков и тепловых карт по сводным таблицам. \\
\hline
\texttt{Web UI} & Графический интерфейс пользователя на Streamlit (FR-10): форма выбора корпуса, векторного хранилища, LLM по ролям, профиля защиты, типа атаки, числа сидов; запуск \texttt{Experiment Orchestrator}; отображение таблицы метрик и предоставление артефактов прогона для скачивания. Является тонкой надстройкой над \texttt{Experiment Orchestrator} и не дублирует его логику. \\
\hline
\end{tabular}
\end{table}

\paragraph{Роли LLM.} Внутри комплекса языковые модели разделяются на четыре роли (таблица~\ref{tab:llm_roles}). Физически за двумя или более ролями может стоять одна и та же модель, но интерфейсы и параметры вызова разделяются программно для прозрачности результатов и упрощения замены отдельных моделей.

\begin{table}[H]
\centering
\caption{Роли LLM в программном комплексе и рекомендуемые модели.}
\label{tab:llm_roles}
\begin{tabular}{|p{0.13\textwidth}|p{0.26\textwidth}|p{0.26\textwidth}|p{0.25\textwidth}|}
\hline
\textbf{Роль} & \textbf{Использование} & \textbf{Модель / провайдер по умолчанию} & \textbf{Режим вызова} \\
\hline
\texttt{GENERATOR} & RAG-генерация ответа $G$; Draft-генерация в стадии 3 HARD & Llama~3.1~8B Instruct через Ollama (локально, дёшево, основной массовый вызов) либо Gemini~2.0~Flash через API & $T \in [0; 0{,}3]$, $\text{max\_tokens}$ умеренное \\
\hline
\texttt{VERIFIER} & Draft-then-Verify (стадия 3 HARD; профиль \texttt{ragfort}); требуется качественная оценка соответствия draft контексту и наличия утечки & Gemini~2.0~Flash через OpenAI-compat API (умная модель, доступ через API) & $T = 0$, короткий ответ, парсинг скаляра $V \in [0,1]$ \\
\hline
\texttt{JUDGE} & LLM-as-a-judge для $\mathbb{I}_{leak}$ и для оценки качества ответов; требуется надёжный детектор утечки & Gemini~2.0~Flash через OpenAI-compat API (или сильнее: Gemini~2.0~Pro/Anthropic Claude Sonnet) & $T = 0$, ответ в формате \texttt{JSON} \\
\hline
\texttt{SEGMENTER} & Декомпозиция запроса $\phi(q)$ для IRD (стадия 1 HARD); задача NLP-сегментации, допустима более лёгкая модель & Gemini~2.0~Flash через API; для отладки~--- Llama~3.1~8B через Ollama & $T = 0$, ответ в формате \texttt{JSON}-массив строк \\
\hline
\end{tabular}
\end{table}

\noindent Принципиально, что роли \texttt{VERIFIER} и \texttt{JUDGE} назначаются на <<умную>> модель через API, а массовый \texttt{GENERATOR} может быть локальной Llama~3.1~8B: это удешевляет основной поток генерации, не жертвуя качеством ключевых оценочных вызовов. Роль \texttt{JUDGE} физически независима от защищающего пайплайна и используется только для измерения метрик (ASR, качество ответов).

\subsubsection{Матрица трассируемости требований}

Для верификации полноты архитектуры приведём явную матрицу соответствия функциональных требований и модулей (таблица~\ref{tab:traceability}). Каждое требование покрывается хотя бы одним модулем, и каждый модуль участвует в реализации хотя бы одного требования.

\begin{table}[H]
\centering
\caption{Матрица трассируемости требований и модулей. Знак $\bullet$~--- модуль участвует в реализации требования.}
\label{tab:traceability}
\small
\begin{tabular}{|l|c|c|c|c|c|c|c|c|c|c|}
\hline
& \rotatebox{75}{\texttt{Corpus Mgr}}
& \rotatebox{75}{\texttt{Index Bld}}
& \rotatebox{75}{\texttt{Poisoned-Idx}}
& \rotatebox{75}{\texttt{RAG Pipeline}}
& \rotatebox{75}{\texttt{Attack Suite}}
& \rotatebox{75}{\texttt{Defense Mgr}}
& \rotatebox{75}{\texttt{HARD Cascade}}
& \rotatebox{75}{\texttt{Metrics Eng}}
& \rotatebox{75}{\texttt{Orchestrator}}
& \rotatebox{75}{\texttt{Web UI}} \\
\hline
FR-1  & $\bullet$ & $\bullet$ &           &           &           &           &           &           &           & $\bullet$ \\
\hline
FR-2  &           &           &           & $\bullet$ &           &           &           & $\bullet$ & $\bullet$ & $\bullet$ \\
\hline
FR-3  &           &           & $\bullet$ &           & $\bullet$ &           &           &           &           & $\bullet$ \\
\hline
FR-4  &           &           &           & $\bullet$ & $\bullet$ &           &           & $\bullet$ & $\bullet$ &           \\
\hline
FR-5  &           &           &           & $\bullet$ &           & $\bullet$ & $\bullet$ &           &           & $\bullet$ \\
\hline
FR-6  &           &           &           & $\bullet$ & $\bullet$ & $\bullet$ & $\bullet$ & $\bullet$ & $\bullet$ &           \\
\hline
FR-7  &           &           &           &           &           &           & $\bullet$ &           & $\bullet$ &           \\
\hline
FR-8  &           &           &           &           &           &           & $\bullet$ & $\bullet$ &           &           \\
\hline
FR-9  &           &           &           &           &           &           &           & $\bullet$ & $\bullet$ & $\bullet$ \\
\hline
FR-10 & $\bullet$ &           &           &           &           & $\bullet$ &           &           & $\bullet$ & $\bullet$ \\
\hline
\end{tabular}
\end{table}

\subsection{Дизайн эксперимента}
\label{sec:experiment_design}

\subsubsection{Факторный план}

Эксперимент строится как полный факторный план по двум измерениям: \textbf{тип атаки} ($a \in \{\text{PI}, \text{Backdoor}, \text{SECRET}\}$) и \textbf{профиль защиты} ($p \in \{\texttt{none}, \texttt{basic-filters}, \texttt{ragfort}, \texttt{hard}\}$), что даёт $3 \times 4 = 12$ сочетаний (attack, defense). Дополнительно для каждого профиля защиты выполняется отдельный benign-прогон ($a = \varnothing$) на $\mathcal{Q}_{legit}$ для измерения FPR, BPD и Latency Overhead. Профили защит и атак сведены в таблицах~\ref{tab:defense_profiles} и~\ref{tab:attack_profiles}.

\begin{table}[H]
\centering
\caption{Профили защиты в эксперименте.}
\label{tab:defense_profiles}
\begin{tabular}{|p{0.20\textwidth}|p{0.68\textwidth}|}
\hline
\textbf{Профиль} & \textbf{Содержимое} \\
\hline
\texttt{none} & Базовая RAG-система без защит. Используется для измерения исходного качества и исходной уязвимости к атакам. \\
\hline
\texttt{basic-filters} & Простые защитные меры: входной \texttt{InputFilter}, фильтрация контекста \texttt{DataFilter} и robust system prompt~--- эквивалент уровня курсовой работы. \\
\hline
\texttt{ragfort} & \emph{Полная} black-box-адаптация архитектуры RAGFort~\cite{RAGFort2025} (\S\ref{sec:def_sota}): Draft-then-Verify (адаптация Constrained Cascade Generation) $+$ Topic-Consistent Re-ranking (адаптация Contrastive Re-indexing). Главный baseline для оценки вклада оригинальных компонентов HARD. \\
\hline
\texttt{hard} & Полный предлагаемый метод (\S\ref{sec:def_hard}): \texttt{ragfort} (DtV $+$ TCR) $+$ оригинальные компоненты IRD, Output Leak Scanner и Risk-Budget Thresholding. Сравнение с профилем \texttt{ragfort} напрямую измеряет вклад этих трёх компонентов. \\
\hline
\end{tabular}
\end{table}

\begin{table}[H]
\centering
\caption{Атаки, используемые в эксперименте.}
\label{tab:attack_profiles}
\begin{tabular}{|p{0.20\textwidth}|p{0.68\textwidth}|}
\hline
\textbf{Атака} & \textbf{Экспериментальный сценарий} \\
\hline
Prompt Injection & Атакующий контролирует только пользовательский запрос и добавляет вредоносную инъекцию $\delta$ к легитимной части $q_{benign}$; индекс не модифицируется. \\
\hline
Backdoor & Атакующий внедряет отравленный документ в корпус через один из штатных каналов ингеста (открытые источники, пользовательский upload, RSS/коннекторы; см.~\S\ref{sec:threat_model}). Эксперимент моделирует результат этого внедрения сборкой отдельного poisoned-индекса $\mathcal{I}_{pois}$ на корпусе с добавленным отравленным документом и не делает различий по конкретному каналу попадания. Прямое изменение готового индекса находится вне модели угроз и не моделируется. \\
\hline
SECRET & Атакующий запрос строится как $x = I_{ext} \oplus O_{jail} \oplus T_{retr}$. $O_{jail}$ оптимизируется LLM-as-an-Optimizer, $T_{retr}$ генерируется с учётом кластерной структуры корпуса (Cluster-Focused Triggering). \\
\hline
\end{tabular}
\end{table}

\subsubsection{Процедура одного прогона}

Один полный прогон эксперимента состоит из следующих шагов:
\begin{enumerate}[label=\arabic*.]
    \item Подготовить корпус $\mathcal{D}$, набор легитимных вопросов $\mathcal{Q}_{legit}$ и набор секретов $S_{priv}$.
    \item Построить чистый индекс $\mathcal{I}$ и измерить baseline-качество RAG-системы без атак и защит.
    \item Построить poisoned-индекс $\mathcal{I}_{pois}$, добавив отравленный документ в корпус и переиндексировав его.
    \item Откалибровать пороги стадии 1 HARD: вычислить $s_{\text{IRD}}(q)$ на бенигн-выборке $\mathcal{Q}_{train}^{ben}$ и установить $\theta_1^{(0)}$ как $0{,}95$-квантиль распределения.
    \item Для каждого профиля $p \in \{\texttt{none}, \texttt{basic-filters}, \texttt{ragfort}, \texttt{hard}\}$ выполнить benign-прогон на $\mathcal{Q}_{legit}$ и вычислить FPR, BPD, $\Delta T$.
    \item Для каждой пары $(a, p)$ выполнить прогон на $\mathcal{Q}_{atk}^{a}$ и вычислить ASR. Для Backdoor-атак использовать $\mathcal{I}_{pois}$, для остальных~--- $\mathcal{I}$.
    \item Для пар с $p = \texttt{hard}$ дополнительно сохранить трассировки стадий и вычислить StagePassRate, StageCorrelation, $\mathrm{AUC}_{\text{IRD}}$, $\kappa$.
    \item Сформировать сводную таблицу результатов и графики.
\end{enumerate}

Шаги 5--7 повторяются для $S$ значений seed. По умолчанию $S = 1$ (seed $= 42$): этого достаточно для отладочных и итеративных прогонов и обеспечивает выполнение бюджета времени NFR-1. Для финального прогона перед защитой работы рекомендуется $S = 3$ (seed $\in \{41, 42, 43\}$); в этом случае итоговые метрики приводятся в виде $\mathrm{mean} \pm \mathrm{std}$. FPR измеряется на наборе $\mathcal{Q}_{legit}$, в который входят как вопросы из $\mathcal{Q}_{golden}$ ($n = 30$ в полном режиме), так и часть бенигн-тренинга, не использованная для калибровки порогов.

\subsubsection{Параметры по умолчанию}

Параметры, фиксируемые на старте эксперимента, приведены в таблице~\ref{tab:experiment_params}. Эти значения соответствуют полному режиму \texttt{full}; режимы \texttt{fast} и \texttt{smoke} получаются из них уменьшением размеров выборок (см. примечание под таблицей). Режим \texttt{smoke} специально предназначен для проверки работоспособности UI и пайплайна за минуты, а не за час.

\begin{table}[H]
\centering
\caption{Параметры эксперимента по умолчанию (режим \texttt{full}).}
\label{tab:experiment_params}
\begin{tabular}{|p{0.30\textwidth}|p{0.20\textwidth}|p{0.40\textwidth}|}
\hline
\textbf{Группа} & \textbf{Параметр} & \textbf{Значение} \\
\hline
Корпус и индекс  & размер чистого корпуса         & 200--500 документов (Wikipedia) + 3 базовых .md \\
\cline{2-3}
                 & размер чанка / overlap         & 512 токенов / 64 \\
\cline{2-3}
                 & top-$k$ retrieval               & 5 \\
\cline{2-3}
                 & candidate-$k$                   & 20 \\
\hline
Выборки          & $|\mathcal{Q}_{golden}|$        & 30 \\
\cline{2-3}
                 & $|\mathcal{Q}_{train}^{ben}|$   & 50 \\
\cline{2-3}
                 & $|\mathcal{Q}_{atk}^{\text{PI}}|$ & 30 \\
\cline{2-3}
                 & $|\mathcal{Q}_{atk}^{\text{BD}}|$ & 30 \\
\cline{2-3}
                 & $|\mathcal{Q}_{atk}^{\text{SECRET}}|$ & 30 (с 5--8 итерациями оптимизации на вопрос) \\
\cline{2-3}
                 & seeds (по умолчанию)            & $\{42\}$ ($S = 1$); для финала $\{41, 42, 43\}$ ($S = 3$) \\
\hline
HARD/IRD         & $k_{retr}$ (top-$k$ для отпечатка фрагмента) & 5 \\
\cline{2-3}
                 & макс. число фрагментов $m$      & 5 \\
\cline{2-3}
                 & $\theta_1^{(0)}$                 & $0{,}95$-квантиль $s_{\text{IRD}}$ на $\mathcal{Q}_{train}^{ben}$ \\
\hline
HARD/TCR         & $K_{top}$ (число кластеров k-means) & 12 \\
\cline{2-3}
                 & $\alpha, \beta, \gamma$         & $0{,}6;\ 0{,}3;\ 0{,}1$ \\
\cline{2-3}
                 & $\theta_2^{(0)}$                & $0{,}5$ \\
\hline
HARD/Verifier    & $\theta_3^{(0)}$ (для $1 - V$)   & $0{,}5$ \\
\hline
HARD/Scanner     & $\tau_{copy}$ (LCSR-порог)      & $0{,}7$ \\
\cline{2-3}
                 & $\theta_4^{(0)}$                & $0{,}5$ \\
\hline
Risk-Budget      & $\lambda_1, \lambda_2, \lambda_3, \lambda_4$ & $0{,}0;\ 0{,}3;\ 0{,}3;\ 0{,}3$ \\
\hline
LLM              & \texttt{GENERATOR}              & Llama 3.1 8B Instruct через Ollama (массовый вызов, локально); $T = 0{,}2$ \\
\cline{2-3}
                 & \texttt{VERIFIER, JUDGE} & Gemini 2.0 Flash через OpenAI-compat API (умная модель); $T = 0$ \\
\cline{2-3}
                 & \texttt{SEGMENTER}             & Gemini 2.0 Flash через API; $T = 0$ \\
\hline
\end{tabular}
\end{table}

\noindent \textbf{Режим \texttt{fast}} (флаг \texttt{--fast} CLI) применяет следующие сокращения: $|\mathcal{Q}_{golden}| = 10$, $|\mathcal{Q}_{atk}| = 10$ для каждого типа атаки, число итераций оптимизации SECRET на вопрос~--- 3, $S = 1$ (один seed). Используется для отладки и регрессионных проверок.

\textbf{Режим \texttt{smoke}} (флаг \texttt{--smoke}) дополнительно сокращает выборки до $|\mathcal{Q}_{golden}| = 5$, $|\mathcal{Q}_{atk}| = 5$ на каждый тип атаки, выполняет 1 итерацию SECRET и пропускает Backdoor (опционально), завершая полный цикл для всех профилей за единицы минут. Режим предназначен исключительно для смоук-тестов UI, CLI-конфигурации и подключений к внешним LLM-провайдерам, и не предназначен для содержательных выводов.

\subsubsection{Артефакты и анализ результатов}

Итоговыми артефактами эксперимента являются:
\begin{itemize}[label=$\bullet$]
    \item \texttt{summary.csv}~--- сводная таблица по всем парам $(a, p)$ для прямой вставки в экспериментальную главу;
    \item \texttt{summary.json}~--- та же таблица в машинно-читаемом формате;
    \item \texttt{report.json}~--- детальный отчёт по каждому отдельному прогону, включая трассировку стадий каскада;
    \item \texttt{experiment.log}~--- журнал событий с временными метками для воспроизводимости;
    \item \texttt{provenance.json}~--- описание окружения и параметров (\S\ref{sec:software_requirements});
    \item графические артефакты: heatmap ASR по парам $(a, p)$, гистограммы FPR/BPD/$\Delta T$ по профилям, ROC-кривая IRD на бенигн против атакующих запросов, stage-pass breakdown для профиля \texttt{hard}, матрица StageCorrelation, кривая $\kappa$.
\end{itemize}

Анализ результатов проводится в трёх плоскостях: (i) сравнение эффективности (ASR при разных профилях) и стоимости (FPR, BPD, $\Delta T$) метода HARD относительно базовых профилей; (ii) декомпозиция вклада стадий каскада HARD по StagePassRate; (iii) эмпирическая проверка допущений теоремы~\ref{thm:cascade_bound} через метрики StageCorrelation и $\kappa$. Такая структура анализа позволяет не только зафиксировать итоговое снижение успешности атак, но и объяснить, за счёт каких компонентов оно достигается, что необходимо для содержательной экспериментальной главы диссертации.

\subsubsection{Запуск экспериментов через интерфейс пользователя}
\label{sec:ui_workflow}

В соответствии с FR-10 все перечисленные выше параметры~--- корпус, векторное хранилище, LLM по ролям, профиль защиты, тип атаки, режим (\texttt{full}/\texttt{fast}/\texttt{smoke}), число сидов~--- доступны для выбора через Web-UI без правки исходного кода. Интерфейс реализуется как единое Streamlit-приложение и состоит из трёх смысловых блоков: (а)~форма конфигурации прогона, (б)~раздел запуска и наблюдения прогресса, (в)~раздел результатов с таблицей метрик и кнопками скачивания артефактов прогона. UI вызывает те же функции \texttt{Experiment Orchestrator}, что и CLI-команда \texttt{attack-rag-experiment}; результаты прогонов, запущенных из UI и из CLI, идентичны при идентичных параметрах. Это обеспечивает воспроизводимость и не вводит отдельного <<UI-only>>-кода обработки.

\newpage
\addcontentsline{toc}{section}{Список литературы}

\begin{thebibliography}{99}

% Статьи по RAG
\bibitem{Lewis2020RAG}
Lewis, P., Perez, E., Piktus, A., Petroni, F., Karpukhin, V., Goyal, N., Küttler, H., Lewis, M., Yih, W., Rocktäschel, T., \& Riedel, S. (2020). Retrieval-Augmented Generation for Knowledge-Intensive NLP Tasks. \textit{arXiv preprint arXiv:2005.11401}.
\url{https://arxiv.org/abs/2005.11401}

\bibitem{Shao2024Survey}
Shao, Y., Yu, H., Li, S., Wang, L., Weng, Y., Zhao, B., Liu, T., Chen, G., Tang, J., Yang, M., Huang, F., Si, L., Yan, Y., Wang, H., Qian, W., \& Wang, K. (2024). A Survey on Retrieval-Augmented Generation for Large Language Models. \textit{arXiv preprint arXiv:2312.10997}.
\url{https://arxiv.org/abs/2312.10997}

% Общие проблемы LLM (обоснование RAG)
\bibitem{Brown2020FewShot}
Brown, T. B., Mann, B., Ryder, N., Subbiah, M., Kaplan, J., Dhariwal, P., Neelakantan, A., Shyam, P., Sastry, G., Askell, A., Agarwal, S., Herbert-Voss, A., Krueger, G., Shevlane, T., Tang, E., Amodei, D., McCandlish, S., Babuschkin, B., \& Sutskever, I. (2020). Language Models are Few-Shot Learners. \textit{arXiv preprint arXiv:2005.14165}.
\url{https://arxiv.org/abs/2005.14165}

% Обзоры угроз LLM/RAG
\bibitem{Qi2024Survey}
Qi, X., Yuan, T., Kang, B. Y., Ye, N., Hu, S., Zhou, C., Pang, T., Wang, Y., Liu, K., Wang, Y., Xu, J., Ji, S., Yang, G., Zhao, X., Xiong, K., Ge, Y., \& Li, Z. (2024). A Survey of Attacks and Defenses for Large Language Models. \textit{arXiv preprint arXiv:2307.10760}.
\url{https://arxiv.org/abs/2307.10760}

% Атаки Prompt Injection
\bibitem{Qi2023PI}
Qi, X., Zeng, Y., Yuan, T., Liu, X., Wang, Y., Yao, M., Zhang, B., Zhu, X., \& Xiong, K. (2023). Follow My Instruction: Evaluating the Robustness of Large Language Models to Prompt Injection Attacks. \textit{arXiv preprint arXiv:2302.12173}.
\url{https://arxiv.org/abs/2302.12173}

% Атаки Data Poisoning / Backdoor
\bibitem{Gu2017Badnets}
Gu, T., Liu, K., Wang, B., Meng, X., Xu, W., Tian, Y., Ji, S., \& Li, Z. (2017). BadNets: Identifying Vulnerabilities in the Machine Learning Model Supply Chain. \textit{arXiv preprint arXiv:1708.06733}.
\url{https://arxiv.org/abs/1708.06733}

\bibitem{Peng2024BackdoorRAG}
Peng, J., Zhang, M., Qi, X., Zeng, Y., Wang, Y., Xiong, K., \& Ji, S. (2024). Backdoor Attacks on Retrieval-Augmented Generation Systems. \textit{arXiv preprint arXiv:2402.11437}.
\url{https://arxiv.org/abs/2402.11437}

% Инструменты и техники
\bibitem{Reimers2019SBERT}
Reimers, N., \& Gurevych, I. (2019). Sentence-BERT: Sentence Embeddings using Siamese BERT-Networks. \textit{arXiv preprint arXiv:1908.10084}.
\url{https://arxiv.org/abs/1908.10084}

% Подходы к защите
\bibitem{Zhang2024RobustRAG}
Zhang, C., Han, S., Zhang, Q., Guo, H., Zhai, Y., Liu, J., Wang, B., Chang, Y., Zhang, X., \& Yang, J. (2024). RobustRAG: A Framework for Enhancing Retrieval Augmented Generation. \textit{arXiv preprint arXiv:2403.10221}.
\url{https://arxiv.org/abs/2403.10221}

% LLM Model Used
\bibitem{distilgpt2}
Sanh, V., Debut, L., Chaumond, J., \& Wolf, T. (2019). DistilBERT, a distilled version of BERT: smaller, faster, cheaper and lighter. \textit{arXiv preprint arXiv:1910.01108}.
\url{https://arxiv.org/abs/1910.01108}

% Hugging Face Transformers Library
\bibitem{transformers}
Wolf, T., Debut, L., Sanh, V., Chaumond, J., Delangue, C., Moi, A., \dots \& Rush, A. M. (2020). Transformers: State-of-the-Art Natural Language Processing. \textit{arXiv preprint arXiv:1910.03771}.
\url{https://arxiv.org/abs/1910.03771}

% SOTA 2025
\bibitem{SECRET2025}
He, Y., Chen, Y., Li, Y., Shao, S., Qi, L., Li, B., Tao, D., \& Qin, Z. (2025). External Data Extraction Attacks against Retrieval-Augmented Large Language Models. \textit{arXiv preprint arXiv:2510.02964}.
\url{https://arxiv.org/abs/2510.02964}

\bibitem{LeakSealer2025}
Panebianco, F., Bonfanti, S., Trovò, F., \& Carminati, M. (2025). LeakSealer: A Semisupervised Defense for LLMs Against Prompt Injection and Leakage Attacks. \textit{arXiv preprint arXiv:2508.00602}.
\url{https://arxiv.org/abs/2508.00602}

\bibitem{ControlNET2025}
Yao, H., Shi, H., Chen, Y., Jiang, Y., Wang, C., \& Qin, Z. (2025). ControlNET: A Firewall for RAG-based LLM Systems. \textit{arXiv preprint arXiv:2504.09593}.
\url{https://arxiv.org/abs/2504.09593}

\bibitem{RAGFort2025}
Li, Q., Pan, M., Xiong, K., Su, G., Shen, Z., Liu, Y., Sun, B., Peng, H., \& Zhang, X. (2025). RAGFort: Dual-Path Defense Against Proprietary Knowledge Base Extraction in Retrieval-Augmented Generation. \textit{arXiv preprint arXiv:2511.10128}.
\url{https://arxiv.org/abs/2511.10128}

\end{thebibliography}

\end{document}
